\documentclass[11pt]{amsart}
\usepackage{amsmath,amssymb,amsthm}
\usepackage[margin=1.1in]{geometry}
\usepackage{hyperref}

\newtheorem{theorem}{Theorem}
\newtheorem{lemma}[theorem]{Lemma}
\newtheorem{proposition}[theorem]{Proposition}
\newtheorem{corollary}[theorem]{Corollary}
\newtheorem{conjecture}[theorem]{Conjecture}
\newtheorem{question}[theorem]{Question}
\theoremstyle{definition}
\newtheorem{definition}[theorem]{Definition}
\newtheorem{remark}[theorem]{Remark}
\newtheorem{example}[theorem]{Example}

\newcommand{\cF}{\mathcal F}
\newcommand{\Ev}{\mathrm{Ev}}
\newcommand{\Od}{\mathrm{Od}}

\title{Endpoint sets of path decompositions at the Gallai budget}
\author{Justin Williams}
\date{November 2026 (working draft)}



\begin{document}
\begin{abstract}
For a graph $G$ of even order $n$ with even-degree vertices $\Ev$, call $S\subseteq\Ev$
\emph{realisable} if some partition of $E(G)$ into simple paths has exactly two ends at every vertex
of $S$, none at the other even vertices, and exactly one at every odd vertex. Such a partition has
$|\Od|/2+|S|$ paths, so Gallai's conjecture (in its rooted even-order form) asks that in every graph
with $p(G)=n/2$ every even vertex lies in a realisable set of size $|\Ev|/2$.
We record the elementary structure of the realisable family: it is an up-set, it contains $\Ev$ and
$\Ev\setminus\{u\}$ for every even $u$ (both from Lov\'asz's theorem on an apex extension), and it is
characterised by the parity-tightness of the graph with private leaves attached at $S$.
We show by explicit examples that three natural strengthenings fail: a complementary-pair property,
and uniformity of the realisable half-sets in $2$-connected graphs, both with and without the
hypothesis $p(G)=n/2$. Exhaustive computation through order six (and order eight for the stated
cases) supports the conjecture that in every graph with $p(G)=n/2$ the minimal realisable sets are
the bases of a loopless matroid. This conjecture would settle all deletion degrees of the
three-connected rooted test simultaneously.
We prove that the first two steps of the corresponding descent always succeed: in a connected
graph of even order with at least four even vertices, for all distinct even $y,w$ some set
$\Ev\setminus\{w,v\}$ with $v\ne y$ is realisable. The proof uses only Lov\'asz's theorem and Fan's
addible-edge transfers. For four even vertices this is the whole conjecture: the realisable pairs
of a connected critical graph are the bases of a loopless rank-two matroid, and every even vertex
lies in at least two of them. All results of this section are formalised in Lean~4.

A second route, through \emph{Fan descent}, controls which pair is removed. The step
$\Ev\setminus\{u\}\to\Ev\setminus\{u,v\}$ is equivalent to the parity bound $p=|\Od|/2$ for a
graph with exactly two even vertices, and we prove it whenever $u$ and $v$ are not twins in the
subgraph induced by the even vertices. Without connectivity, faithful descent therefore takes its
second step except in the two configurations of Corollary~\ref{cor:fd2strong}; for connected
graphs Theorem~\ref{thm:two} closes both (Remark~\ref{rem:twostep}).
Iterating the argument gives a calculus of edge-deletion moves (Fan certificates). It yields a
sufficient condition for the rooted test that is uniform in the deletion degree. The main
statements are formalized in Lean.

Finally (Section~\ref{sec:H}) we test whether a stronger induction hypothesis can rescue the
branching case.
\begin{itemize}
\item One prescribed vertex reduces the case of two even neighbours. So Gallai's conjecture follows
from the branching case with at least three even neighbours.
\item An oracle that knows every realisable end profile of $G-e$ certifies the rooted test, through
Fan's moves on $e$, at every even root of every connected graph of order at most eight, with one
exception: $K_8$ minus a perfect matching.
\item A counting barrier explains the exception. For $K_{2m}$ minus a perfect matching, $m\ge4$, no
Fan transfer and no short new path can be the last move of a certificate, whatever the seed. Its
endpoint matroid is uniform, so the obstruction lies in the moves, not in the hypothesis.
\end{itemize}
The positive counterpart is the \emph{flower move}. When the receipts at a hub stop at Fan's
equality, the paths close into cycles through the hub. A path that takes one edge from each cycle
then recombines them into one more path, with two ends at the hub. For $K_{2m}$ minus a perfect
matching, both the failure of Fan's test and the rooted test at every vertex are already Theorem~10.21
of~\cite{SP}, which opens Walecki's Hamilton cycles along a path. Our contributions are the move in
general: equality forces a flower, and the dichotomy ``receipt or flower''. As an application, the
order-$10$ graph $B'$ of~\cite{MM} satisfies the rooted test at every vertex.
The endpoint matroid conjecture holds for all $1{,}082$ connected critical graphs of order eight.
\end{abstract}
\maketitle

\section{Setting}

All graphs are finite and simple. A \emph{path} is a nonempty simple path; a \emph{partition} of a
graph is a partition of its edge set into paths, and $p(G)$ is the least size of a partition. For a
partition $P$ and a vertex $v$ write $b_P(v)$ for the number of ends of paths of $P$ at $v$, counted
with multiplicity. Since every path through $v$ uses one or two edges at $v$ according as $v$ is an
end or an interior vertex of it,
\begin{equation}\label{eq:parity}
b_P(v)\equiv d(v)\pmod 2,\qquad \sum_v b_P(v)=2|P|.
\end{equation}
Write $\Od=\Od(G)$ and $\Ev=\Ev(G)$ for the sets of odd-degree and of even positive-degree
vertices, $o=|\Od|$, $e=|\Ev|$. For a graph without isolated vertices $o+e=n$, and $o$ is even.

Gallai's conjecture asserts $p(G)\le\lceil n/2\rceil$ for connected $G$. It is equivalent to the
following rooted statement at even order~\cite{REC,SP}: for every connected $H$ of order $2k$ and
every $r\in V(H)$, the graph $H^r$ obtained by attaching a new leaf at $r$ satisfies $p(H^r)\le k$.
If $p(H)<k$ this holds trivially (split a path at an interior visit to $r$, or extend a path already
ending at $r$), and if $d(r)$ is odd a minimum partition already has an end at $r$. Hence the rooted
statement is equivalent to:

\begin{quote}
\textbf{(R)} If $H$ is connected of even order $n$ with $p(H)=n/2$, then every even vertex is an end
of some partition of $H$ into $n/2$ paths.
\end{quote}
We call a graph of even order with $p(H)=n/2$ \emph{critical}.

\section{Realisable sets}

\begin{definition}
A set $S\subseteq\Ev$ is \emph{realisable} if $G$ has a partition $P$ with $b_P=1$ on $\Od$, $b_P=2$
on $S$ and $b_P=0$ on $\Ev\setminus S$. Such a partition has $o/2+|S|$ paths by \eqref{eq:parity}.
Write $\cF(G)$ for the family of realisable sets.
\end{definition}

\begin{lemma}[Up-closure]\label{lem:up}
If $S\in\cF(G)$ and $S\subseteq T\subseteq \Ev$ then $T\in\cF(G)$.
\end{lemma}

\begin{proof}
It suffices to add one vertex $v\in\Ev\setminus S$. In a partition realising $S$ the vertex $v$ has no
ends and positive degree, so it is interior to some path; splitting that path at $v$ gives two
nonempty paths and two ends at $v$, and changes no other end count.
\end{proof}

We use Lov\'asz's theorem~\cite{L68}: \emph{a graph on $N$ vertices with at most one vertex of even
degree has a partition into $\lfloor N/2\rfloor$ paths.} (Isolated vertices are allowed; apply the
theorem componentwise.)

\begin{lemma}[Faithful partitions]\label{lem:faithful}
$\Ev\in\cF(G)$ for every graph $G$ without isolated vertices.
\end{lemma}

This is Cai's theorem on faithful path partitions~\cite{Cai}; we include the short deduction from
Lov\'asz's theorem since the same device gives Lemma~\ref{lem:minus}.

\begin{proof}
If $\Ev=\emptyset$ the claim is Lov\'asz's theorem itself. Otherwise form $G^+$ by adding a vertex
$z$ adjacent to every vertex of $\Ev$. In $G^+$ every vertex of $G$ has odd degree and $d(z)=e$.
If $e$ is odd, all $N=n+1$ vertices of $G^+$ are odd, so $N$ is even and Lov\'asz gives $N/2$ paths,
each vertex being an end exactly once (there are $N$ odd vertices and $N$ ends). If $e$ is even, $z$
is the unique even vertex, $N=n+1$ is odd, Lov\'asz gives $(N-1)/2$ paths and the $N-1$ odd vertices
absorb all $N-1$ ends, so $z$ has none. In both cases every vertex of $G$ is an end of exactly one path of the partition of
$G^+$. Delete $z$, discarding any piece without edges. Only the neighbours of $z$ change their end
counts, and each $w\in N(z)$ loses exactly its edge $zw$. Let $Q$ be the path carrying $zw$. If $w$ is
interior to $Q$, then $w$ becomes an end of a piece of $Q$, so $w$ now has two ends. If $w$ is an end of
$Q$, then $zw$ is the end edge of $Q$ at $w$; removing it removes the unique end of $w$ (the piece of
$Q$ on the far side of $z$ survives, if nonempty, and the piece $\{w\}$ is discarded), so $w$ now has no
end. Thus every $w\in N(z)$ ends with $0$ or $2$ ends and every other vertex keeps its count. Finally
each $w\in N(z)$ with no end has positive even degree in $G$, so it is interior to some path, and we
split that path at $w$. The resulting partition realises $\Ev$.
\end{proof}

\begin{lemma}[Removing one vertex]\label{lem:minus}
Let $G$ have even order and no isolated vertex. For every $u\in\Ev$, $\Ev\setminus\{u\}\in\cF(G)$.
\end{lemma}

\begin{proof}
Since $n$ is even, $e$ is even. Add a vertex $z$ adjacent to every vertex of $\Ev\setminus\{u\}$. In the
resulting graph $G^+$ on $N=n+1$ vertices the vertices of $\Ev\setminus\{u\}$ and of $\Od$ are odd,
$d(z)=e-1$ is odd, and $u$ is the unique even vertex. Lov\'asz's theorem gives a partition into
$(N-1)/2=n/2$ paths; the $n$ odd vertices of $G^+$ absorb all $n$ ends, so each has exactly one end
and $u$ has none. In particular $z$ is the end of exactly one path. Delete $z$ and repair exactly as
in the proof of Lemma~\ref{lem:faithful}: each $w\in N(z)=\Ev\setminus\{u\}$ ends with $0$ or $2$
ends and is split if it has none, while $u$ and the odd vertices are not adjacent to $z$ and keep $0$
and $1$ ends respectively.
\end{proof}

Lemma~\ref{lem:minus} says that the first step of any descent from the faithful partition succeeds,
for every choice of the vertex kept.

\begin{lemma}[Leaf dictionary]\label{lem:dict}
For $S\subseteq\Ev$ let $G^S$ be $G$ with a new leaf attached at each vertex of $S$. Then $S\in\cF(G)$
if and only if $p(G^S)=|\Od(G^S)|/2$.
\end{lemma}

\begin{proof}
$\Od(G^S)=\Od(G)\cup S\cup L$, where $L$ is the set of new leaves, so $|\Od(G^S)|=o+2|S|$. If $P$
realises $S$, append each new leaf edge to one of the two paths ending at its attachment vertex
(different leaves at different vertices use different ends). This gives $o/2+|S|$ paths. Conversely a
partition of $G^S$ into $|\Od(G^S)|/2$ paths has exactly one end at each odd vertex of $G^S$ and none at
the vertices of $\Ev\setminus S$. A leaf edge $s\ell$ is the end edge of its path. If that path has
another edge, trimming the leaf leaves $s$ with two ends. If the path is the single edge $s\ell$,
deleting it leaves $s$ with no ends, and $s$ is then interior to another path, at which we split. In
every case the result realises $S$.
\end{proof}

\begin{corollary}\label{cor:R}
For a critical graph $H$, every realisable set has size at least $e/2$, and (R) holds for $H$ at the
even vertex $r$ if and only if $r$ lies in a realisable set of size exactly $e/2$, or $r$ is an end in
a partition of $H$ into $n/2$ paths in which some odd vertex has at least three ends.
\end{corollary}

\begin{proof}
A realisable $S$ gives $o/2+|S|\ge p(H)=n/2=(o+e)/2$ paths. The second statement restates (R),
distinguishing whether the partition has a simple profile.
\end{proof}

\section{Three strengthenings that fail}

\begin{proposition}[Complementary pairs]\label{prop:cp}
There are connected graphs $G$ of odd order, outside the odd semi-cliques, with no pair of partitions
$P_1,P_2$ such that $b_{P_1}=b_{P_2}=1$ on $\Od$ and $\{b_{P_1}(v),b_{P_2}(v)\}=\{0,2\}$ for every
$v\in\Ev$. An example is $K_3\vee \overline{K_4}$.
\end{proposition}

\begin{proof}
Let $G=K_3\vee\overline{K_4}$. The three triangle vertices have degree $6$ and the other four have
degree $3$, so $n=7$, $o=4$, $e=3$. Such a pair would satisfy $|P_1|+|P_2|=(2o+2e)/2=7$. Since a
vertex of degree $6$ needs at least three paths, both partitions have at least three paths, and one of
them, say $P_1$, has exactly three paths and exposes exactly $|P_1|-o/2=1$ even vertex. But in a
partition into three paths a vertex of degree $6$ is interior to all three paths, so no even vertex can
be an end of $P_1$, a contradiction. Note $p(G)=3=\lfloor 7/2\rfloor$.
\end{proof}

\begin{proposition}[Non-critical $2$-connected graphs]\label{prop:uhcstar}
Let $G$ have vertices $0,\dots,5$ and edges $01,02,03,05,14,25,35,45$. Then $G$ is $2$-connected, all
degrees are even, $p(G)=2<3$, and the half-set $\{0,1,4\}$ of even vertices is not realisable.
\end{proposition}

\begin{proof}
Suppose a partition has ends exactly at $0,1,4$ (two each) and none at $2,3,5$. Since $2$ and $3$
are interior, the paths contain the segments $0\,2\,5$ and $0\,3\,5$. At the interior vertex $5$ the
four edges $52,53,50,54$ are paired into two transitions. Pairing $52$ with $53$ closes the cycle
$0\,2\,5\,3\,0$; pairing $52$ with $50$ closes $0\,2\,5\,0$; likewise $53$ cannot pair with $50$. So
$50$ pairs with $54$, and then $52$ must pair with $53$, a contradiction. The value $p(G)=2$ is
attained by the paths $2\,0\,1\,4\,5\,3$ and $2\,5\,0\,3$.
\end{proof}

\begin{proposition}[Critical $2$-connected graphs]\label{prop:uhc}
Let $G$ have vertices $0,\dots,7$, with $0$ adjacent to all others, $6$ adjacent to $1,2,7$ and $7$
adjacent to $3,4,5,6$. Then $G$ is $2$-connected, $p(G)=4=n/2$ (forced by $d(0)=7$), the even vertices
are $1,\dots,6$, and neither half-set $\{3,4,5\}$ nor $\{1,2,6\}$ is realisable. All other $18$
half-sets are realisable, and they form the bases of a sparse paving matroid.
\end{proposition}

\begin{proof}[Proof of non-realisability]
For $\{3,4,5\}$: the even vertices $1,2,6$ would be interior, giving segments $0\,1\,6$ and $0\,2\,6$.
At $6$ (degree $4$, interior) pairing $61$ with $62$, $61$ with $60$, or $62$ with $60$ closes a
cycle through $0$, so both $61$ and $62$ must pair with $67$, which is impossible. For $\{1,2,6\}$:
$3,4,5$ are interior, giving segments $0\,3\,7$, $0\,4\,7$, $0\,5\,7$. At $7$, which has degree $5$ and
exactly one end, no two of $73,74,75$ can be paired with each other, nor with $70$, without closing a
cycle through $0$. So at most one of them pairs with $76$ and at most one is the end edge, leaving the
third unmatched, a contradiction.
The realisability of the other $18$ half-sets was verified by exact search (Section~\ref{app:comp}).
\end{proof}

\begin{proposition}[Critical graphs need not be forced]\label{prop:notforced}
Let $1\le j<k$. In $K_{2k+1}$ choose a matching $\{a_ib_i:1\le i\le j\}$, delete it, and add a new vertex
$w$ adjacent to $a_1,\dots,a_j$. The resulting graph $G$ has order $n=2k+2$, maximum degree $n-2$,
exactly $(n/2-1)(n-1)$ edges and at most $j+1$ odd vertices, and $p(G)\ge n/2$. It is $2$-connected for
$j\ge 2$ and $j$-connected for $j\ge3$.
\end{proposition}

\begin{proof}
$|E(G)|=\binom{2k+1}{2}-j+j=k(2k+1)=(n/2-1)(n-1)$. A path has at most $n-1$ edges, so a partition
into $k$ paths consists of $k$ Hamilton paths. Each contains $w$ and uses at least one of its $j<k$
edges, and distinct paths use distinct edges; this is impossible. The degrees are $2k=n-2$ at the
untouched vertices and at the $a_i$, $2k-1$ at the $b_i$, and $j$ at $w$, so the odd vertices are the
$b_i$ together with $w$ when $j$ is odd. $K_{2k+1}$ minus a matching is $(2k-1)$-connected, and $w$ has $j$
neighbours in it, which gives the connectivity claims.
\end{proof}

For $j=2$, $k=3$ this is the graph $K_7-\{ac,bd\}$ plus $w\sim\{a,b\}$. For $j=3$, $k=4$ it is a
$3$-connected critical graph on ten vertices that is not forced by any trivial bound. Its six even
vertices have a uniform endpoint matroid: all $20$ half-sets are realisable, as verified by exact
search.

So the trivial lower bounds $o/2$, $\lceil\Delta/2\rceil$, $\lceil m/(n-1)\rceil$ do not detect all
critical graphs, and (R) has content beyond the dense regime.

\section{The endpoint matroid conjecture}

\begin{conjecture}[Endpoint matroid]\label{conj:emc}
If $H$ is connected and critical, the minimal members of $\cF(H)$ are the bases of a matroid on
$\Ev(H)$ of rank $e/2$ without loops.
\end{conjecture}

By Lemma~\ref{lem:up}, the conjecture says that $\cF(H)$ is the family of spanning sets of a loopless
matroid. Looplessness is exactly (R), via the simple-profile half of Corollary~\ref{cor:R}. Hence the
conjecture implies Gallai's conjecture.

\begin{conjecture}[Faithful descent]\label{conj:fd}
If $G$ is connected of even order, then every $S\in\cF(G)$ with $|S|>e/2$ and every $r\in S$ admit
some $u\in S\setminus\{r\}$ with $S\setminus\{u\}\in\cF(G)$.
\end{conjecture}

Starting from Lemma~\ref{lem:faithful} and descending, Conjecture~\ref{conj:fd} yields (R). By
Lemma~\ref{lem:minus} its first step always succeeds, and by Theorem~\ref{thm:two} of
Section~\ref{sec:four} so does its second: the conjecture holds for every $S$ with $|S|\ge e-1$.

\begin{proposition}[Small even sets]\label{prop:small}
Let $H$ be connected and critical. If $e\le 2$, every even vertex lies in a realisable set of size $e/2$.
If $e=4$, the same holds provided the floor bound holds for odd-order graphs with a leaf and exactly
three even vertices (which follows from the theorem of Botler and Sambinelli~\cite{BS}).
\end{proposition}

\begin{proof}
For $e=2$ and $\Ev=\{r,s\}$, Lemma~\ref{lem:minus} with $u=s$ shows that $\{r\}$ is realisable. For $e=4$,
attach a leaf at $r$. The resulting odd-order graph has exactly the three even vertices $\Ev\setminus\{r\}$,
and its even subgraph has maximum degree at most two. By~\cite[Theorem~13]{BS} it satisfies the floor
bound unless its even subgraph is a triangle and every odd vertex has at least two even neighbours.
The new leaf has no even neighbour, so this exception does not occur, and the graph has a partition
into $\lfloor (n+1)/2\rfloor=n/2$ paths. Trimming the leaf exposes $r$ twice, or releases a path that
pays for a split at $r$.
\end{proof}

Thus looplessness in Conjecture~\ref{conj:emc} is known for $e\le4$. The first case with content in the
matroid part is $e=4$: if $\{1,2\}$ and $\{3,4\}$ are realisable, the conjecture requires $\{1,3\}$ or
$\{1,4\}$ to be realisable as well. Section~\ref{sec:four} proves the conjecture for $e=4$
(Theorem~\ref{thm:emc4}), and Corollary~\ref{cor:rooted4} reproves the case $e=4$ of
Proposition~\ref{prop:small} without~\cite{BS}.

\begin{remark}[Relation to the deletion tests]
In the three-connected reduction of~\cite{MR}, $J$ has a unique even vertex $z$ of degree $d$,
$B=J-z$ has $\Ev(B)=N_J(z)$, and one asks for a partition of $B$ into $(|J|-1)/2$ paths with two ends
at a prescribed $r\in N_J(z)$. Lov\'asz's theorem on $J$ restricted to $B$ is exactly the faithful
partition of Lemma~\ref{lem:faithful}, and the test asks for $d/2$ descent steps keeping $r$. Both
conjectures treat every $d$ at once.
\end{remark}

\section{Two descent steps, and four even vertices}\label{sec:four}

In this section every graph has no isolated vertices. For a partition $D$ of a graph $K$ and a
vertex $u$ let $Z_D(u)=\{v\in N_K(u): b_D(v)=0\}$ be the set of \emph{passing neighbours} of $u$.
If $D$ realises $S$, then a vertex has no ends exactly when it is an even vertex outside $S$, so
\begin{equation}\label{eq:passing}
Z_D(u)=N_K(u)\cap(\Ev(K)\setminus S).
\end{equation}
The only tool besides Lov\'asz's theorem is Fan's addible-edge lemma~\cite{Fan}, in its single and
grouped forms (see also~\cite[Lemmas~6 and~7]{BS}).

\begin{lemma}[Fan's transfers]\label{lem:fan}
Let $D$ be a partition of a graph $K$, let $u$ be a vertex, and let $T$ be a set of vertices with
$u\notin T$ and no vertex of $T$ adjacent to $u$. If $\sum_{t\in T}b_D(t)>|Z_D(u)|$, then for some
$t\in T$ the graph $K+ut$ has a partition $D'$ with $|D'|=|D|$, $b_{D'}(u)=b_D(u)+1$,
$b_{D'}(t)=b_D(t)-1$, and $b_{D'}(v)=b_D(v)$ for every other vertex $v$.
\end{lemma}

For $T=\{t\}$ we call this a \emph{single transfer of $ut$ towards $u$}; for larger $T$ a
\emph{grouped transfer}. Both forms are proved in full in the Lean library (\texttt{fan\_single},
\texttt{fan\_grouped}). Deleting one edge $ab$ flips the parities of $a$ and $b$ and of no other
vertex, and this is the only bookkeeping in the proofs below.

\begin{lemma}[Transfer across an even edge]\label{lem:edge}
Let $G$ have even order, and let $s,t,u'$ be distinct even vertices with $st\in E(G)$ and
$su'\notin E(G)$. Then $\Ev\setminus\{t,u'\}\in\cF(G)$.
\end{lemma}

\begin{proof}
Let $K=G-st$. In $K$ the vertices $s$ and $t$ are odd, of degree at least one, and all other
degrees are unchanged, so $K$ has even order, no isolated vertex and $\Ev(K)=\Ev\setminus\{s,t\}$.
Lemma~\ref{lem:minus} applied to $K$ and $u'$ gives a partition $D$ of $K$ realising
$\Ev\setminus\{s,t,u'\}$. By~\eqref{eq:passing} and $su'\notin E(G)$ we get $Z_D(s)=\emptyset$, while
$b_D(t)=1$. A single transfer of $st$ towards $s$ gives a partition of $G$ of the same size in which
$s$ has two ends, $t$ has none, and every other vertex keeps its ends. It realises
$\Ev\setminus\{t,u'\}$.
\end{proof}

\begin{lemma}[Star lemma]\label{lem:star}
Let $y,w,x,z$ be distinct even vertices of $G$ with $y$ adjacent to $w$, $x$ and $z$. Then
$\Ev\setminus\{w,x\}$ or $\Ev\setminus\{w,z\}$ is in $\cF(G)$.
\end{lemma}

\begin{proof}
Let $K=G-\{yw,yx,yz\}$. In $K$ the vertex $y$ has degree $d(y)-3$ and $w,x,z$ have degree one less
than in $G$, so all four are odd and $\Ev(K)=\Ev\setminus\{y,w,x,z\}$. By Lemma~\ref{lem:faithful},
$K$ has a partition $D_0$ realising $\Ev(K)$; no vertex is without ends in $D_0$.
\begin{enumerate}
\item $Z_{D_0}(y)=\emptyset$ and $b_{D_0}(w)=1$. A single transfer of $yw$ towards $y$ gives a
partition $D_1$ of $K+yw$ with $b_{D_1}(y)=2$, $b_{D_1}(w)=0$, and all other end counts unchanged.
\item In $D_1$ the only vertex without ends is $w$, so $|Z_{D_1}(y)|\le1<2=b_{D_1}(x)+b_{D_1}(z)$.
A grouped transfer with $T=\{x,z\}$ gives, for some $t'\in\{x,z\}$, a partition $D_2$ of
$K+yw+yt'$ with $b_{D_2}(y)=3$ and $b_{D_2}(t')=0$.
\item Let $t''$ be the other vertex of $\{x,z\}$. In $D_2$ only $w$ and $t'$ have no ends, so
$|Z_{D_2}(t'')|\le2<3=b_{D_2}(y)$. A single transfer of $yt''$ towards $t''$ gives a partition $D_3$
of $G$ with $b_{D_3}(t'')=2$ and $b_{D_3}(y)=2$.
\end{enumerate}
In $D_3$ the vertices of $\Ev(K)\cup\{y,t''\}$ have two ends, $w$ and $t'$ have none, and every odd
vertex of $G$ has one. So $D_3$ realises $\Ev\setminus\{w,t'\}$.
\end{proof}

\begin{lemma}[Lifting along an odd neighbour]\label{lem:lift}
Let $y$ be an even vertex of $G$, let $t$ be an odd neighbour of $y$, and let $K=G-yt$. In $K$ the
vertex $y$ is odd and $t$ is even, and all other degrees are as in $G$. If $S\in\cF(K)$ contains $t$
and $y$ has no neighbour in $\Ev(K)\setminus S$, then $(S\setminus\{t\})\cup\{y\}\in\cF(G)$.
\end{lemma}

\begin{proof}
Let $D$ realise $S$ in $K$. By~\eqref{eq:passing} and the hypothesis, $Z_D(y)=\emptyset$, while
$b_D(t)=2$. A single transfer of $yt$ towards $y$ raises the ends at $y$ from one to two and lowers
those at $t$ from two to one, which is the profile of $(S\setminus\{t\})\cup\{y\}$ in $G$.
\end{proof}

\begin{theorem}[Two descent steps]\label{thm:two}
Let $G$ be connected of even order with $e\ge4$ even vertices. For all distinct even vertices $y$
and $w$ there is an even vertex $v\notin\{y,w\}$ with $\Ev\setminus\{w,v\}\in\cF(G)$.
\end{theorem}

\begin{proof}
We prove the statement for every graph $G$ of even order, without isolated vertices and with
$e\ge4$, in which some path $P$ joins $y$ to another even vertex, by induction on the length of $P$.

\emph{Case 1: $y$ has an even neighbour.} If $y\sim w$ and $y\not\sim v$ for some even
$v\notin\{y,w\}$, Lemma~\ref{lem:edge} with $(s,t,u')=(y,w,v)$ gives $\Ev\setminus\{w,v\}$. If $y$
is adjacent to every other even vertex, choose distinct even $x,z\notin\{y,w\}$ (as $e\ge4$) and
apply Lemma~\ref{lem:star}. If $y\not\sim w$, then $y\sim t$ for some even $t\notin\{y,w\}$, and
Lemma~\ref{lem:edge} with $(s,t,u')=(y,t,w)$ gives $\Ev\setminus\{t,w\}$; take $v=t$.

\emph{Case 2: $y$ has no even neighbour.} Then the second vertex $t$ of $P$ is odd. Let $K=G-yt$.
It has even order and no isolated vertex ($y$ had degree at least two, and $t$ keeps the next edge
of $P$), and $\Ev(K)=(\Ev\setminus\{y\})\cup\{t\}$ has $e$ elements. The path $P-y$ joins $t$ to an
even vertex of $K$ and is shorter than $P$, so by induction, applied to $K$ with $t$ in place of $y$,
there is $v\in\Ev(K)\setminus\{t,w\}$ with $S=\Ev(K)\setminus\{w,v\}\in\cF(K)$. Now $t\in S$, and
$\Ev(K)\setminus S=\{w,v\}$ consists of even vertices of $G$, which are not adjacent to $y$. By
Lemma~\ref{lem:lift}, $(S\setminus\{t\})\cup\{y\}=\Ev\setminus\{w,v\}$ is in $\cF(G)$.
\end{proof}

Lemma~\ref{lem:minus} is the first step of faithful descent for every choice of the deleted vertex,
and Theorem~\ref{thm:two} is the second step for every choice of the kept vertex $y$ and of the
first deleted vertex $w$. Criticality is not used. Connectivity is needed: for the disjoint union
of a path on three vertices (whose middle vertex is $y$) and $K_5-e$, the even vertices are $y$ and
the three vertices of degree four of $K_5-e$, and no realisable pair contains $y$, because
$p(K_5-e)=3$ forbids realising a single vertex of $K_5-e$.

\begin{corollary}[Four even vertices]\label{cor:four}
Let $G$ be connected of even order with exactly four even vertices. For distinct even $y,x,z$, the
pair $\{y,x\}$ or the pair $\{y,z\}$ is realisable. Equivalently, every even vertex lies in at least
two realisable pairs, and no two non-realisable pairs of even vertices share a vertex.
\end{corollary}

\begin{proof}
Let $w$ be the fourth even vertex, and apply Theorem~\ref{thm:two} to $y$ and $w$: the set
$\Ev\setminus\{w,v\}$ with $v\in\{x,z\}$ is $\{y,z\}$ or $\{y,x\}$.
\end{proof}

In other words, the \emph{partner graph} on $\Ev$, joining two even vertices when they form a
realisable pair, has minimum degree at least two. Its complement is therefore a matching, and the
partner graph is $C_4$, $K_4-e$ or $K_4$.

\begin{theorem}[The endpoint matroid for four even vertices]\label{thm:emc4}
Let $H$ be connected and critical with exactly four even vertices. Then the minimal realisable sets
are exactly the realisable pairs, and they are the bases of a loopless matroid of rank two on
$\Ev(H)$. Its non-trivial parallel classes are the non-realisable pairs, and there are at most two
of them. In particular Conjecture~\ref{conj:emc} holds when $e=4$.
\end{theorem}

\begin{proof}
By Corollary~\ref{cor:R} every realisable set has at least two elements. If $S$ is realisable and
$|S|\ge3$, pick distinct $y,x,z\in S$; by Corollary~\ref{cor:four} $S$ contains the realisable pair
$\{y,x\}$ or $\{y,z\}$. Hence the minimal realisable sets are the realisable pairs. They satisfy
basis exchange: let $B_1=\{a,c\}$ and $B_2$ be realisable pairs with $a\notin B_2$. If $c\in B_2$,
then $B_2=\{b,c\}$ with $b\notin B_1$ and $B_1-a+b=B_2$. If $c\notin B_2=\{b_1,b_2\}$,
Corollary~\ref{cor:four} applied to $c,b_1,b_2$ makes $\{c,b_1\}$ or $\{c,b_2\}$ realisable, and
this is $B_1-a+b_i$. So the realisable pairs are the bases of a matroid of rank two. Every even
vertex lies in a realisable pair (Corollary~\ref{cor:four}), so there are no loops. Two elements
are parallel exactly when they form a non-realisable pair, and by Corollary~\ref{cor:four} these
pairs are disjoint.
\end{proof}

Conversely, since $\Ev\setminus\{u\}$ is realisable for every $u$ (Lemma~\ref{lem:minus}),
Conjecture~\ref{conj:emc} for $e=4$ requires every $3$-set to contain a basis. So it is
equivalent to Corollary~\ref{cor:four} for critical graphs, which is stronger than the statement
that non-realisability of pairs is transitive. Both non-trivial parallel classes can occur. Let $G$ consist of two triangles $a\,x_1x_2$ and
$b\,z_1z_2$ joined by the edge $ab$, so $n=6$ and $\Ev=\{x_1,x_2,z_1,z_2\}$. If neither vertex of
degree two in a triangle is an end, a single path contains all three edges of that triangle, which is
impossible; so each triangle
contains an even vertex with two ends, every partition has at least $(1+1+2+2)/2=3$ paths, and $G$
is critical. The same argument shows that $\{x_1,x_2\}$ and $\{z_1,z_2\}$ are not realisable, while
for instance $\{x_1,z_1\}$ is realised by the paths $x_1x_2abz_2z_1$, $ax_1$ and $bz_1$.

\begin{corollary}[The rooted test for four even vertices]\label{cor:rooted4}
Let $G$ be connected of even order $n$ with exactly four even vertices. Then for every vertex $r$,
the graph $G^r$ obtained by attaching a leaf at $r$ satisfies $p(G^r)\le n/2$. Consequently every
connected graph $L$ of odd order with a pendant leaf and exactly three even vertices satisfies
$p(L)\le\lfloor|L|/2\rfloor$.
\end{corollary}

\begin{proof}
If $r$ is even, Theorem~\ref{thm:two} gives a realisable pair containing $r$; if $r$ is odd, take
any realisable pair. A realising partition has $o/2+2=n/2$ paths and an end at $r$, which the leaf
edge extends. For the second statement let $\ell$ be the leaf of $L$, $s$ its neighbour, and
$G=L-\ell$, a connected graph of even order. If $s$ is odd in $L$, then $\Ev(G)=\Ev(L)\cup\{s\}$ has
four elements and $L=G^s$. If $s$ is even in $L$, then $s$ is odd in $G$ and $G$ has exactly two even
vertices, $\Ev(G)=\{a,b\}$; by Lemma~\ref{lem:minus} the set $\{a\}$ is realisable in $G$, by a
partition into $o(G)/2+1=|G|/2$ paths with an end at $s$, which the leaf edge extends.
\end{proof}

Corollary~\ref{cor:rooted4} is the floor bound $\mathrm{BP}_3$ of the project roadmap. It also follows
from~\cite[Theorem~13]{BS}; the point here is that it needs only Lov\'asz's theorem and Fan's
transfers.

\begin{remark}[Six even vertices]
For $e=6$, Theorem~\ref{thm:two} gives, for every even $r$ and every other even $w$, a realisable
$4$-set containing $r$ and avoiding $w$. The simple-profile form of (R) asks for a realisable
$3$-set containing $r$, that is, one further descent step; by the remark on the deletion tests this
is the open case $d=6$ of the three-connected test. Iterating the proof meets a specific obstruction.
For a third step one may delete an edge $yt$ with $t$ even and apply Theorem~\ref{thm:two} to
$K=G-yt$, obtaining $\Ev(K)\setminus\{w',v'\}\in\cF(K)$; a single transfer of $yt$ towards $y$ then
gives $\Ev\setminus\{t,w',v'\}\in\cF(G)$, provided $K$ is connected and neither $w'$ nor $v'$ is
adjacent to $y$. The vertex $w'$ can be prescribed, but $v'$ cannot. So the third step needs a version
of Theorem~\ref{thm:two} in which the second deleted vertex avoids a prescribed set of vertices. When
$y$ has exactly one even neighbour this obstruction disappears (Theorem~\ref{thm:unique} below).
\end{remark}

\subsection*{$k$-step descent}

For an even vertex $y$, a set $W_0\subseteq\Ev\setminus\{y\}$ and $k\ge1$, we say that $G$ has
\emph{$k$-step descent at $(y,W_0)$} if some $k$-set $R\subseteq\Ev\setminus\{y\}$ with $W_0\subseteq R$
has $\Ev\setminus R\in\cF(G)$. Lemma~\ref{lem:minus} is $1$-step descent at $(y,\{w\})$, and
Theorem~\ref{thm:two} is $2$-step descent at $(y,\{w\})$. For a critical graph with $e=2k$, $k$-step
descent at $(y,\emptyset)$ is the simple-profile form of (R) at $y$. Prescribing $w$ is a convenience
of the induction: at the bottom level $w$ may be a coloop of the endpoint matroid.

\begin{lemma}[General transfer across an even edge]\label{lem:transfer}
Let $s,t$ be adjacent even vertices of $G$ and $K=G-st$. If $S\in\cF(K)$ and $s$ has no neighbour in
$\Ev(K)\setminus S$, then $S\cup\{s\}\in\cF(G)$.
\end{lemma}

\begin{proof}
In $K$ both $s$ and $t$ are odd. Let $D$ realise $S$ in $K$. By~\eqref{eq:passing}, $Z_D(s)=\emptyset$,
and $b_D(t)=1$. A single transfer of $st$ towards $s$ gives $s$ two ends and $t$ none.
\end{proof}

Lemma~\ref{lem:edge} is the case $S=\Ev\setminus\{s,t,u'\}$, supplied by Lemma~\ref{lem:minus}.

\begin{theorem}[Connected reduction]\label{thm:reduction}
Fix $k\ge1$, $m$ and a finite set $W_0$. Suppose that in every connected graph with at least $m$ even
vertices, $k$-step descent holds at $(y,W_0)$ for every even $y\notin W_0$ (with $W_0\subseteq\Ev$)
that has an even neighbour or all of whose edges are bridges. Then it holds for every even
$y\notin W_0$ in every connected graph with at least $m$ even vertices containing $W_0$.
\end{theorem}

\begin{proof}
Induction on $|E(G)|$. If $y$ is not covered by the hypothesis, then $y$ has no even neighbour and some
edge $yt$ is not a bridge, so $t$ is odd and $K=G-yt$ is connected. In $K$ the vertex $y$ is odd and
$t$ is even, so $\Ev(K)=(\Ev\setminus\{y\})\cup\{t\}$ has as many elements as $\Ev$, and
$W_0\subseteq\Ev(K)\setminus\{t\}$. By induction there is a $k$-set $R\subseteq\Ev(K)\setminus\{t\}$
with $W_0\subseteq R$ and $\Ev(K)\setminus R\in\cF(K)$. The vertices of $R$ are even in $G$ and
different from $y$, so they are not adjacent to $y$, and Lemma~\ref{lem:lift} gives
$\Ev\setminus R\in\cF(G)$.
\end{proof}

Neither parity nor criticality is used, and the induction never changes the vertex set, so the
theorem also holds for graphs on a fixed vertex set, in particular of fixed even order; this is the
form combined with Corollary~\ref{cor:three} below (and the form of \texttt{desc\_reduction}).
Connectivity matters for $k\ge3$. The proof of
Theorem~\ref{thm:two} is a path induction that never needs connectivity of the graph, but the
corresponding base statement fails for disconnected graphs when $k=3$. In $(K_5-e)\sqcup(K_5-e)$
($n=10$, $e=6$) every even vertex has an even neighbour, yet every realisable set has at least four
elements, since each copy needs two ($p(K_5-e)=3$). When every edge at $y$ is a bridge, deleting one
of them splits $G$, possibly into parts of odd order, where descent is the floor problem; this case
is not reduced here.

\begin{theorem}[One even neighbour]\label{thm:unique}
Let $k\ge2$, let $G$ be connected with at least $2k+2$ even vertices, let the even vertex $y$ have
exactly one even neighbour $t$, and let $w\ne y$ be even. Suppose that $k$-step descent holds at
$(y',\{w'\})$ in every graph on $V(G)$ without isolated vertices and with at least $2k$ even vertices, for all
distinct even $y',w'$ such that $y'$ is joined by a path to another even vertex. Then $G$ has
$(k+1)$-step descent at $(y,\{w\})$.
\end{theorem}

\begin{proof}
Let $K=G-yt$. Then $y$ and $t$ are odd in $K$, $\Ev(K)=\Ev\setminus\{y,t\}$ has at least $2k\ge4$
elements, and $K$ has no isolated vertex. Since $G$ is connected, every vertex of $K$ is joined in $K$
to $y$ or to $t$; as $\Ev(K)$ has at least four elements, two distinct even vertices $a,b$ of $K$ lie
in one component of $K$. Let $y'$ be
one of them with $y'\ne w$, and let $w'=w$ if $w\ne t$ and $w'$ be the other of $a,b$ if $w=t$. The
hypothesis gives a $k$-set $W\subseteq\Ev(K)\setminus\{y'\}$ with $w'\in W$ and
$\Ev(K)\setminus W\in\cF(K)$. As $t$ is the only even neighbour of $y$ and $W\subseteq\Ev\setminus\{y,t\}$,
the vertex $y$ has no neighbour in $W$, and Lemma~\ref{lem:transfer} gives
$(\Ev(K)\setminus W)\cup\{y\}=\Ev\setminus(W\cup\{t\})\in\cF(G)$. The set $R=W\cup\{t\}$ has $k+1$
elements, contains $w$ and avoids $y$.
\end{proof}

The proof of Theorem~\ref{thm:two} is an induction along a path from $y$ to another even vertex that
uses neither connectivity nor criticality, so it supplies the hypothesis of Theorem~\ref{thm:unique}
for $k=2$, in graphs of even order.

\begin{corollary}[Three steps]\label{cor:three}
Let $G$ be connected of even order with $e\ge6$, and let the even vertex $y$ have exactly one even
neighbour. Then for every even $w\ne y$ some $3$-set $R\ni w$ of even vertices avoiding $y$ has
$\Ev\setminus R\in\cF(G)$. If $e=6$, then $y$ lies in a realisable $3$-set and the rooted test (R)
holds at $y$.
\end{corollary}

In the language of the deletion tests: if $J$ is $3$-connected with unique even vertex $z$ of degree
$6$, then test~(i) holds at every $r\in N(z)$ that has exactly one neighbour in $N(z)$ (apply
Corollary~\ref{cor:three} to $B=J-z$, which is connected, has even order and $\Ev(B)=N(z)$). Combined with
Theorem~\ref{thm:reduction}, three-step descent in connected graphs of even order with $e\ge6$
reduces to two cases: $y$ has at least two even neighbours, or $y$ has no even neighbour and every
edge at $y$ is a bridge.

\subsection*{Bridges, and reduction to the branching case}

For $j\ge1$ let $\mathrm{D}_j$ be the statement: \emph{every connected graph of even order with at
least $2j$ even vertices has $j$-step descent at $(y,\{w\})$ for all distinct even $y,w$}. Let
$\mathrm{B}_j$ be the same statement restricted to vertices $y$ with at least two even neighbours.
$\mathrm{D}_1$ is Lemma~\ref{lem:minus}, and $\mathrm{D}_2$ is Theorem~\ref{thm:two}. For a bridge
$yt$ of $G$ let $A$ and $B$ be the vertex sets of the components of $G-yt$ containing $y$ and $t$, and
write $A^+=G[A\cup\{t\}]$ and $B^+=G[B\cup\{y\}]$. Every neighbour of $y$ other than $t$ lies in $A$,
so $A^+$ is $G[A]$ with the pendant edge $yt$, $y$ has the same degree in $A^+$ as in $G$, and
$\Ev(A^+)=\Ev\cap A$; symmetrically $\Ev(B^+)=\Ev\cap B$.

\begin{lemma}[Bridge gluing]\label{lem:bridge}
Let $yt$ be a bridge of $G$.
\begin{enumerate}
\item If $S_A\in\cF(A^+)$ and $S_B\in\cF(B^+)$, then $S_A\cup S_B\in\cF(G)$.
\item Let $y$ be even in $G$ and $d(t)\ge2$ (so that $G[B]$ has no isolated vertex),
$S_A\in\cF(G[A])$ and $S_B\in\cF(G[B])$, where $t\in S_B$ if $t$ is
even in $G[B]$. Then $S_A\cup(S_B\setminus\{t\})\cup\{y\}\in\cF(G)$.
\end{enumerate}
\end{lemma}

\begin{proof}
(1) Let $D_A,D_B$ realise $S_A,S_B$. The leaf $t$ of $A^+$ is the end of the path $Q_A=P_A\,yt$ of
$D_A$ through $yt$, and the leaf $y$ of $B^+$ is the end of the path $Q_B=yt\,P_B$ of $D_B$; here
$P_A$ ends at $y$ and lies in $G[A]$, $P_B$ starts at $t$ and lies in $G[B]$, and either may be
trivial. Replace $Q_A,Q_B$ by the path $P_A\,yt\,P_B$. The vertex $y$ is an end of the new path
exactly when it was an end of $Q_A$, and $t$ exactly when it was an end of $Q_B$; every other vertex
keeps its ends, and the two leaf ends disappear. So $y$ has $b_{D_A}(y)$ ends, $t$ has $b_{D_B}(t)$
ends, and the result realises $S_A\cup S_B$.
(2) Let $D_A,D_B$ realise $S_A,S_B$. As $y$ is odd in $G[A]$ it has one end in $D_A$, and $t$ is an
end of some path $P$ of $D_B$ (one end if $t$ is odd in $G[B]$, two if $t\in S_B$). Extend $P$ by
the edge $ty$, which is possible as $y\notin B$. Now $y$ has two ends and $t$ one fewer.
\end{proof}

\begin{theorem}[Reduction to the branching case]\label{thm:branching}
For every $k\ge3$, if $\mathrm{B}_3,\dots,\mathrm{B}_k$ hold, then $\mathrm{D}_k$ holds.
\end{theorem}

\begin{proof}
Induction on $k$, starting from $\mathrm{D}_2$. Assume $\mathrm{D}_1,\dots,\mathrm{D}_{k}$ and
$\mathrm{B}_{k+1}$, and prove $(k+1)$-step descent at $(y,\{w\})$ in a connected graph $G$ of even
order with $e\ge2k+2$, by induction on $|E(G)|$. In each case below the pieces we descend in are
connected of even order with fewer edges than $G$. A piece in which $r\ge1$ vertices are to be removed
has at least $2r$ even vertices, so we may apply $\mathrm{D}_r$ (by the outer hypothesis if $r\le k$, by
the inner induction if $r=k+1$), with the prescribed vertex $w$ if $w$ lies in that piece and an
arbitrary one otherwise; a piece with $r=0$ uses Lemma~\ref{lem:faithful}. The numbers $r$ are chosen
below so that they fit the capacities (half the number of even vertices of each piece), add up
correctly, and are positive on the piece containing $w$; this is possible because the capacities add
up to at least the required total and the piece containing $w$ has capacity at least one.

\emph{Case 1: $y$ has at least two even neighbours.} This is $\mathrm{B}_{k+1}$.

\emph{Case 2: $y$ has a bridge $yt$ with $d(t)\ge2$.} Both sides have edges, since $d(y),d(t)\ge2$.
If $|A|$ is odd, then $A^+$ and $B^+$ have even order and $|\Ev(A^+)|+|\Ev(B^+)|=e$. Take
$(r_A)$-step descent in $A^+$ keeping $y$ and $(r_B)$-step descent in $B^+$ with $r_A+r_B=k+1$, and
glue by Lemma~\ref{lem:bridge}(1). If $|A|$ is even, then $G[A]$ and $G[B]$ have even order, $y$ is odd
in $G[A]$, and $\Ev(G[A])=\Ev\cap A\setminus\{y\}$. If $t$ is even in $G$, it is odd in $G[B]$,
$|\Ev(G[A])|+|\Ev(G[B])|=e-2$, and descents with $r_A+r_B=k$ glue by Lemma~\ref{lem:bridge}(2) to
remove the $k$ vertices and $t$. If $t$ is odd in $G$, it is even in $G[B]$,
$|\Ev(G[A])|+|\Ev(G[B])|=e$, and descents with $r_A+r_B=k+1$, keeping $t$ in $G[B]$, glue by
Lemma~\ref{lem:bridge}(2).

\emph{Case 3: $y$ has exactly one even neighbour $t$, and $yt$ is not a bridge.} Then $K=G-yt$ is
connected of even order with $e-2\ge2k$ even vertices, and the proof of Theorem~\ref{thm:unique} applies
with $\mathrm{D}_k$ in $K$ (no component argument is needed).

\emph{Case 4: otherwise.} Then $y$ has no even neighbour and every bridge at $y$ is pendant. Not every
edge at $y$ is pendant, for otherwise $G$ would be a star with one even vertex. So some edge $yt$ is
neither pendant nor a bridge, $t$ is odd, and the induction step in the proof of
Theorem~\ref{thm:reduction} applies.
\end{proof}

Since $\mathrm{B}_k$ is a special case of $\mathrm{D}_k$, Theorem~\ref{thm:branching} says that
$\mathrm{D}_k$ for all $k$ is equivalent to $\mathrm{B}_k$ for all $k\ge3$. For a critical graph with
$e=2k$, $\mathrm{D}_k$ places every even vertex in a realisable set of size $e/2$, which is (R). Hence:

\begin{corollary}\label{cor:branching}
If $\mathrm{B}_k$ holds for every $k\ge3$, then Gallai's conjecture holds. If $\mathrm{B}_3$ holds,
then the rooted test (R) holds for every connected graph with six even vertices.
\end{corollary}

The proof of Theorem~\ref{thm:branching} works verbatim without the prescribed vertex $w$, and Case~2
may be applied before Case~1. So the hypothesis can be weakened to $\mathrm{B}^0_k$: \emph{in every
connected graph of even order with at least $2k$ even vertices, every even vertex $y$ with at least two
even neighbours and no bridge to a vertex of degree at least two lies in a realisable set
$\Ev\setminus R$ with $|R|=k$}. If $\mathrm{B}^0_k$ holds for all $k\ge3$, Gallai's conjecture holds
(Lean: \texttt{BranchBase}, \texttt{gallai\_of\_branchBase}). The statement $\mathrm{B}_3$ (three-step
descent at an even vertex with at least two even neighbours) is the current target of the third
descent step. Theorem~\ref{thm:branching} also disposes of the
all-bridges case left open in Theorem~\ref{thm:reduction}.

\begin{remark}[Formalisation]\sloppy
Lemmas~\ref{lem:edge}--\ref{lem:lift}, Theorem~\ref{thm:two}, Corollaries~\ref{cor:four}
and~\ref{cor:rooted4} and Theorem~\ref{thm:emc4} are formalised in Lean~4 with Mathlib, in the file
\texttt{lean/GallaiLean/EMCFour.lean} of the project repository (\texttt{realisable\_of\_edge},
\texttt{realisable\_of\_star}, \texttt{realisable\_lift}, \texttt{two\_step\_descent},
\texttt{emc\_four}, \texttt{endpoint\_matroid\_four}, \texttt{rooted\_four}; the second statement of
Corollary~\ref{cor:rooted4} is \texttt{rooted\_le\_four}, the rooted test for at most four even
vertices). The matroid of
Theorem~\ref{thm:emc4} is built as a Mathlib \texttt{Matroid} and proved loopless. There is no
\texttt{sorry}, and the only axioms are \texttt{propext}, \texttt{Classical.choice} and
\texttt{Quot.sound}. Lov\'asz's theorem is an explicit hypothesis (\texttt{LovaszStatement}); Fan's
transfers are proved in the library. The $k$-step material is in \texttt{lean/GallaiLean/KStep.lean}:
Lemma~\ref{lem:transfer} (\texttt{realisable\_transfer}), Theorem~\ref{thm:reduction}
(\texttt{desc\_reduction}), Theorem~\ref{thm:unique} (\texttt{desc\_succ\_of\_unique}),
Corollary~\ref{cor:three} (\texttt{desc\_three\_of\_unique}) and the reduction of three-step descent
to the two residual cases (\texttt{desc\_three\_reduction}). Primed versions of the last two discharge
the Lov\'asz hypothesis with the library's proof of Lov\'asz's theorem. Lemma~\ref{lem:bridge},
Theorem~\ref{thm:branching} and Corollary~\ref{cor:branching} are in
\texttt{lean/GallaiLean/Bridge.lean} (\texttt{glue\_bridge}, \texttt{glue\_append},
\texttt{dstmt\_of\_branching}, \texttt{rooted\_six\_of\_bstmt3}, \texttt{gallai\_of\_branching}),
with no hypotheses at all: the statements $\mathrm{D}_j$ and $\mathrm{B}_j$ are quantified over all
vertex types, and pieces are handled as subgraphs on the same vertex set and transported through the
induced graph on their non-isolated vertices.
\end{remark}

\section{Fan descent, Fan certificates, and the second step for arbitrary pairs}\label{sec:second}

Throughout this section $G$ has no isolated vertex. By Lemma~\ref{lem:minus} every set
$\Ev\setminus\{u\}$ is realisable. We show that the next step, from $\Ev\setminus\{u\}$ to
$\Ev\setminus\{u,v\}$, fails only between vertices that are twins in the subgraph induced by the
even vertices, and we give a calculus of moves that iterates the argument. The mechanism answers
the question of how Fan's lemma, which never changes the number of paths, can drive a descent,
which lowers it by one at each step: the merge is paid for by deleting an edge between two even
vertices, which makes both odd; Lov\'asz's theorem is applied to the smaller graph and Fan's lemma
puts the edge back.

\begin{remark}[Lemmas~\ref{lem:faithful} and~\ref{lem:minus} at every order]\label{rem:anyorder}
The even-order hypothesis of Lemma~\ref{lem:minus} is not needed. For $u\in\Ev$ the graph
$G^{\Ev\setminus\{u\}}$ has the single even vertex $u$ and $n+e-1$ vertices; since $n\equiv o+e\equiv
e\pmod 2$, this number is odd, so Lov\'asz's theorem gives a partition into
$(n+e-2)/2=|\Od(G^{\Ev\setminus\{u\}})|/2$ paths, and Lemma~\ref{lem:dict} shows that
$\Ev\setminus\{u\}$ is realisable. In the same way $G^{\Ev}$ has no even vertex, which gives
Lemma~\ref{lem:faithful}. (The Lean proofs use the apex form and keep the hypothesis $n$ even. The main statements of this
section are formalized for even $n$; not formalized are the moves (S) and (C) of
Proposition~\ref{prop:moves}, the classification in Corollary~\ref{cor:fd2},
Corollary~\ref{cor:fd2strong}, the assembly of Theorem~\ref{thm:b}, and (O2), (O3) and the
cut-vertex criterion of Section~\ref{sec:obstructions}.)
\end{remark}

\subsection{Tight graphs}

\begin{definition}
A graph $H$ is \emph{tight} if it has a partition into $|\Od(H)|/2$ paths. By~\eqref{eq:parity} this
is the same as a partition with exactly one end at every odd vertex and none at the even vertices,
that is, $\emptyset\in\cF(H)$.
\end{definition}

By Lemma~\ref{lem:dict}, $S\in\cF(G)$ if and only if $G^S$ is tight. For $S=\Ev\setminus\{u,v\}$ the
graph $G^S$ has exactly two even vertices, $u$ and $v$. Hence the second descent step is the
\emph{two-even-vertex Lov\'asz problem}: which graphs with exactly two even vertices meet the
parity bound $p=|\Od|/2$? For graphs with at most one even vertex Lov\'asz's theorem says that all
of them do. (The apex graph $G+z+\ell$ of the project notes, with $z$ joined to $\Ev\setminus\{u,v\}$ and a
leaf $\ell$ at $z$, only gives a sufficient condition; $G^{\Ev\setminus\{u,v\}}$ gives an
equivalence.)

\subsection{Fan's lemma}

\begin{lemma}[Fan~\cite{Fan}, Lemma~3.4]\label{lem:fan-b}
Let $ab\in E(G)$ and let $D'$ be a partition of $G-ab$. If fewer than $b_{D'}(b)$ neighbours $x$ of
$a$ in $G-ab$ have $b_{D'}(x)=0$, then $G$ has a partition $D$ with $|D|=|D'|$,
$b_D(a)=b_{D'}(a)+1$, $b_D(b)=b_{D'}(b)-1$, and $b_D(x)=b_{D'}(x)$ for every $x\ne a,b$.
\end{lemma}

The proof is Fan's chain argument: a path ending at $b$ either avoids $a$, and then absorbs the edge
$ba$, or reaches $a$ through an edge $ya$, which is traded for $ba$ after reversing the segment
from $b$ to $y$, moving the missing end to $y$. A counting argument shows that a chain of such
reroutes reaches a vertex with an end on a path avoiding $a$. The lemma, together with the grouped
version, is formalized without restriction (\texttt{fan\_single}, \texttt{fan\_grouped} in
\texttt{lean/GallaiLean/Fan.lean}).

\subsection{Fan descent and Fan transfer}

\begin{lemma}[Fan descent]\label{lem:fandescent}
Let $S\subseteq\Ev$, let $u\ne w$ be vertices of $S$ with $uw\in E(G)$, and suppose that no
neighbour of $w$ lies in $\Ev\setminus S$. If $S\setminus\{u,w\}$ is realisable in $G-uw$, then
$S\setminus\{u\}$ is realisable in $G$.
\end{lemma}

\begin{proof}
Put $G'=G-uw$. The vertices $u$ and $w$ have odd degree in $G'$ and all other degrees are unchanged,
so $\Ev(G')=\Ev\setminus\{u,w\}$ and $\Od(G')=\Od\cup\{u,w\}$. Let $D'$ realise
$S'=S\setminus\{u,w\}$ in $G'$. Then $b_{D'}=1$ on $\Od\cup\{u,w\}$, $b_{D'}=2$ on $S'$, and
$b_{D'}=0$ exactly on $\Ev\setminus S$. A neighbour $x$ of $w$ in $G'$ is neither $w$ nor $u$, so if
$b_{D'}(x)=0$ then $x$ is an even vertex of $G$ outside $S$; by hypothesis there is no such
neighbour. As $b_{D'}(u)=1>0$, Lemma~\ref{lem:fan-b} with $(a,b)=(w,u)$ gives a partition $D$ of $G$
with $|D|=|D'|$, $b_D(w)=2$, $b_D(u)=0$ and all other end counts as in $D'$. Since $u$ and $w$ are
even in $G$, $D$ realises $S\setminus\{u\}$.
\end{proof}

The sizes match as they must: $|D'|=|\Od(G')|/2+|S'|=(o+2)/2+|S|-2=o/2+|S\setminus\{u\}|$.

\begin{lemma}[Fan transfer]\label{lem:fantransfer}
Let $S\subseteq\Ev$, $u\in\Ev\setminus S$ and $w\in\Od$ with $uw\in E(G)$, and suppose every even
neighbour $x\ne u$ of $w$ lies in $S$. If $S$ is realisable in $G-uw$, then $S$ is realisable in
$G$.
\end{lemma}

\begin{proof}
In $G'=G-uw$ the vertex $u$ is odd and $w$ is even, so $S\subseteq\Ev(G')$ and $w\notin S$. Let $D'$
realise $S$ in $G'$; then $b_{D'}(u)=1$ and $b_{D'}(w)=0$. A neighbour $x$ of $w$ in $G'$ with
$b_{D'}(x)=0$ is an even vertex of $G$ different from $u$ and $w$ and not in $S$; there is none.
Lemma~\ref{lem:fan-b} with $(a,b)=(w,u)$ gives $D$ on $G$ with $b_D(w)=1$, $b_D(u)=0$ and all other
counts unchanged, and $D$ realises $S$ in $G$.
\end{proof}

Fan transfer does not shrink $S$; it moves the role of the excluded even vertex from $u$ to its odd
neighbour $w$, at the cost of the edge $uw$.

\subsection{The second step}

\begin{theorem}[Second descent step]\label{thm:second}
Let $u,v,w$ be distinct even vertices with $uw\in E(G)$ and $vw\notin E(G)$. Then
$\Ev\setminus\{u,v\}$ is realisable.
\end{theorem}

\begin{proof}
The graph $G'=G-uw$ has no isolated vertex, since $u$ and $w$ have positive even degree in $G$, and
$\Ev(G')=\Ev\setminus\{u,w\}$ contains $v$. By Lemma~\ref{lem:minus} (Remark~\ref{rem:anyorder}),
$\Ev(G')\setminus\{v\}=\Ev\setminus\{u,v,w\}$ is realisable in $G'$. Apply
Lemma~\ref{lem:fandescent} with $S=\Ev\setminus\{v\}$: $u,w\in S$, and the only vertex of
$\Ev\setminus S$ is $v$, which is not adjacent to $w$. The conclusion is that
$S\setminus\{u\}=\Ev\setminus\{u,v\}$ is realisable.
\end{proof}

\begin{definition}
Two distinct even vertices $u,v$ are \emph{E-twins} if they have the same neighbours among the
other even vertices, $N(u)\cap(\Ev\setminus\{u,v\})=N(v)\cap(\Ev\setminus\{u,v\})$; that is, they are
twins (true or false) in the E-subgraph $G[\Ev]$.
\end{definition}

\begin{corollary}[Non-realisable pairs are E-twins]\label{cor:twins}
If $\Ev\setminus\{u,v\}$ is not realisable, then $u$ and $v$ are E-twins.
\end{corollary}

\begin{proof}
If some even $w\notin\{u,v\}$ is adjacent to exactly one of $u,v$, apply Theorem~\ref{thm:second}
to $w$ and to the pair $u,v$ ordered so that the first vertex is the neighbour of $w$.
\end{proof}

Being twins in a graph (with or without the edge between the two) is an equivalence relation: if
$u\equiv v$ and $v\equiv w$ are distinct, then a vertex $x\notin\{u,v,w\}$ is adjacent to all three or
to none, and $v\sim u\Leftrightarrow u\sim w\Leftrightarrow w\sim v$, where the outer
equivalences come from $v\equiv w$ and $u\equiv v$. A class with at least three vertices is a
clique or an independent set. So the pairs $\{u,v\}$ with $\Ev\setminus\{u,v\}\notin\cF$ lie inside
the twin classes of $G[\Ev]$, which are modules of $G[\Ev]$. (Not every twin pair fails; for
instance in $K_{2,4}$ plus the edge between the two vertices of degree $5$, every pair of the four
degree-$2$ vertices is realisable.)

\begin{corollary}[Faithful descent, second step]\label{cor:fd2}
Let $e\ge4$, and let $u\ne r$ be even vertices. Some even $v\notin\{u,r\}$ has
$\Ev\setminus\{u,v\}\in\cF(G)$, so Conjecture~\ref{conj:fd} can take its second step from
$\Ev\setminus\{u\}$ keeping $r$, unless $\Ev\setminus\{r\}$ lies in the twin class of $u$. The latter
happens for some $u$ (and then for all $u\in\Ev\setminus\{r\}$) exactly when $G[\Ev]$ is one of
\[
K_e,\qquad \overline{K_e},\qquad K_{e-1}+r\ \text{($r$ isolated in $G[\Ev]$)},\qquad
K_{1,e-1}\ \text{(a star centred at $r$)}.
\]
\end{corollary}

\begin{proof}
The first statement is Corollary~\ref{cor:twins}. If $\Ev\setminus\{r\}$ lies in one twin class
$C$, then $|C|\ge e-1\ge3$, so $G[\Ev\setminus\{r\}]$ is complete or edgeless; and $r$ is adjacent to
all or none of $\Ev\setminus\{r\}$, because any two vertices of $C$ other than $r$ agree on $r$.
Conversely each of the four graphs makes all of $\Ev\setminus\{r\}$ pairwise twins.
\end{proof}

So the second step of faithful descent is proved, for every root, except when the E-subgraph is
complete or edgeless up to the edges at the root; Theorem~\ref{thm:pair} below removes most of
these exceptions, and for connected $G$ Theorem~\ref{thm:two} removes all of them
(Remark~\ref{rem:twostep}).

\subsection{Fan certificates}

Lemmas~\ref{lem:fandescent} and~\ref{lem:fantransfer} compose. For $X\subseteq\Ev$ call a sequence of
moves a \emph{Fan certificate for $X$} if it is empty and $|X|\le1$, or its first move is a pair
$(x,w)$ with $x\in X$, $xw\in E(G)$, and $w$ adjacent to no vertex of $X\setminus\{x\}$, such that
either
\begin{itemize}
\item[(E)] $w$ is even and $w\notin X$, and the rest is a Fan certificate for $X\setminus\{x\}$ in
$G-xw$ (\emph{even step}); or
\item[(O)] $w$ is odd with $d(w)\ge3$, and the rest is a Fan certificate for
$(X\setminus\{x\})\cup\{w\}$ in $G-xw$ (\emph{odd transfer}).
\end{itemize}

\begin{theorem}[Soundness]\label{thm:fancert}
If $X\subseteq\Ev$ has a Fan certificate, then $\Ev\setminus X$ is realisable.
\end{theorem}

\begin{proof}
Induction on the length. For $|X|\le1$ use Lemma~\ref{lem:faithful} or~\ref{lem:minus}. In an even
step, $\Ev(G-xw)=\Ev\setminus\{x,w\}$, the degrees in $G-xw$ stay positive, and by induction
$(\Ev\setminus\{x,w\})\setminus(X\setminus\{x\})=S\setminus\{x,w\}$ is realisable in $G-xw$, where
$S=\Ev\setminus(X\setminus\{x\})$. The neighbours of $w$ in $\Ev\setminus S=X\setminus\{x\}$ are
excluded, so Lemma~\ref{lem:fandescent} gives $S\setminus\{x\}=\Ev\setminus X$. In an odd transfer,
$\Ev(G-xw)=(\Ev\setminus\{x\})\cup\{w\}$ (here $d(w)\ge3$ keeps $w$ non-isolated), so by induction
$\Ev(G-xw)\setminus((X\setminus\{x\})\cup\{w\})=\Ev\setminus X$ is realisable in $G-xw$, and
Lemma~\ref{lem:fantransfer} with $S=\Ev\setminus X$ lifts it to $G$.
\end{proof}

Theorem~\ref{thm:second} is the certificate $((u,w))$ for $X=\{u,v\}$. With even steps only, a
certificate is a list $(x_j,w_j),\dots,(x_2,w_2)$ followed by the single vertex $x_1$, with all
vertices distinct and even, $x_iw_i\in E(G)$, and $w_i$ adjacent to none of $x_1,\dots,x_{i-1}$
(edges deleted earlier do not matter, since they join $x_k$ to $w_k$ with $k>i$). In other words
$\{x_iw_i\}$ is a \emph{semi-induced matching} of $G[\Ev]$ in the ordering
$x_1,\dots,x_j$.

\begin{corollary}[A uniform sufficient condition for (R)]\label{cor:semiinduced}
Let $H$ be critical, $r\in\Ev(H)$ and $j=e/2$. Suppose there are distinct even vertices
$x_1,\dots,x_j\ne r$ and $w_2,\dots,w_j\notin\{x_1,\dots,x_j\}$ with $x_iw_i\in E(H)$ and
$x_kw_i\notin E(H)$ for all $k<i$. Then $\Ev\setminus\{x_1,\dots,x_j\}$ is a realisable set of size
$e/2$ containing $r$, so (R) holds for $H$ at $r$. This applies in particular when $H[\Ev]$
contains an induced matching $\{x_2w_2,\dots,x_jw_j\}$ avoiding $r$ and an even vertex
$x_1\ne r$ outside it with no neighbour among $w_2,\dots,w_j$.
\end{corollary}

In the three-connected reduction, $\Ev(B)=N_J(z)$, so the condition concerns only the subgraph
of $J$ induced by the neighbourhood of the deletion vertex, and it has the same form for every
deletion degree $d$. It fails when $J[N(z)]$ has no edges, which is where odd transfers, the pair
step and the reductions below are needed.

\subsection{Twin classes: the pair step}

Twin classes are handled by Fan's grouped lemma, used in Lean as \texttt{fan\_grouped}: \emph{if
$T$ is a set of non-neighbours of $a$, $a\notin T$, and fewer than $\sum_{t\in T}b_D(t)$ neighbours
of $a$ have no end, then for some $t\in T$ the graph $G+ta$ has a partition $D'$ with $|D'|=|D|$,
$b_{D'}(a)=b_D(a)+1$, $b_{D'}(t)=b_D(t)-1$, and the same counts elsewhere}~\cite[Lemma~3.3]{Fan}.

\begin{theorem}[Pair step]\label{thm:pair}
Let $u,a,a'$ be distinct even vertices and let $w\ne u$ be a common neighbour of $a$ and $a'$.
Suppose that either $w$ is even with $d(w)\ge4$, or $w$ is odd and $aa'\notin E(G)$ or neither $a$
nor $a'$ is adjacent to $u$. Then $\Ev\setminus\{u,a\}$ or $\Ev\setminus\{u,a'\}$ is realisable.
\end{theorem}

\begin{proof}
Put $G_2=G-aw-a'w$. In $G_2$ the vertices $a,a'$ are odd, $w$ keeps its parity, and no vertex is
isolated ($d_{G_2}(w)=d(w)-2\ge1$, since an odd $w$ with two neighbours has $d(w)\ge3$). So
$\Ev(G_2)=\Ev\setminus\{a,a'\}$ contains $u$, and Lemma~\ref{lem:minus} gives a partition $D_2$ of
$G_2$ realising $\Ev\setminus\{u,a,a'\}$: one end at $a$, at $a'$ and at the odd vertices, two at
the other even vertices except $u$, none at $u$. In particular $b_{D_2}(w)=2$ if $w$ is even and
$1$ if $w$ is odd, and the only vertex next to $w$ in $G_2$ without an end is (possibly) $u$.

Apply the grouped lemma at $w$ with $T=\{a,a'\}$: since $b_{D_2}(a)+b_{D_2}(a')=2>1$, for some
$t\in\{a,a'\}$ the graph $G-t'w$, where $t'$ is the other vertex, has a partition $D'$ with
$b_{D'}(t)=0$, $b_{D'}(w)=b_{D_2}(w)+1$ and the same counts elsewhere. Now apply
Lemma~\ref{lem:fan} to the edge $t'w$, moving an end from $w$ to $t'$. The vertices next to $t'$ in
$G-t'w$ without an end are among $t$ and $u$. If $w$ is even, $b_{D'}(w)=3>2$; if $w$ is odd,
$b_{D'}(w)=2$, and the hypothesis says that $t$ and $u$ are not both neighbours of $t'$. Either way
the lemma applies and gives a partition $D$ of $G$ with $b_D(t')=2$, $b_D(w)=b_{D_2}(w)$,
$b_D(t)=b_D(u)=0$ and all other counts as in $D_2$, which realises $\Ev\setminus\{u,t\}$.
\end{proof}

\begin{corollary}[Second step of faithful descent]\label{cor:fd2strong}
Let $e\ge4$, $u\ne r$ even. Some even $v\notin\{u,r\}$ has $\Ev\setminus\{u,v\}\in\cF(G)$ unless
\begin{itemize}
\item[(i)] $G[\Ev]$ has no edges and every odd vertex has at most one neighbour in
$\Ev\setminus\{u,r\}$, or
\item[(ii)] $e=4$ and $G[\Ev]$ is a triangle together with the isolated vertex $r$.
\end{itemize}
\end{corollary}

\begin{proof}
Suppose no $v$ works. By Corollary~\ref{cor:fd2}, $G[\Ev\setminus\{r\}]$ is complete or edgeless
and $r$ is adjacent to all or none of $\Ev\setminus\{r\}$. If $G[\Ev\setminus\{r\}]$ is complete
and $e\ge5$ (odd orders are allowed in this section), any three vertices $a,a',w$ of
$\Ev\setminus\{u,r\}$ satisfy Theorem~\ref{thm:pair} ($d(w)\ge e-2\ge3$, hence $d(w)\ge4$ as $w$
is even). If it is complete, $e=4$ and $r$ is adjacent to $\Ev\setminus\{r\}$, take
$\{a,a'\}=\Ev\setminus\{u,r\}$ and $w=r$, which has degree at least $3$, hence at least $4$. If
$G[\Ev\setminus\{r\}]$ is edgeless and $r$ is adjacent to all of it, take $w=r$ again, now with
$d(r)\ge e-1\ge3$. In the remaining edgeless case, an odd vertex with two neighbours
$a,a'\in\Ev\setminus\{u,r\}$ satisfies Theorem~\ref{thm:pair} since $aa'\notin E$. What is left is
(i), and (ii).
\end{proof}

When the first removed vertex may be chosen as well, which is what goal~(b) of the second step
asks, even the case (i) disappears.

\begin{theorem}\label{thm:b}
Let $G$ be connected of even order with $e\ge4$, and let $r\in\Ev$. Then there are distinct
$u,v\in\Ev\setminus\{r\}$ with $\Ev\setminus\{u,v\}$ realisable, unless $e=4$ and $G[\Ev]$ is a
triangle together with the isolated vertex $r$.
\end{theorem}

\begin{proof}
By Corollary~\ref{cor:fd2strong} only the case that $G[\Ev]$ has no edges and every odd vertex has
at most one neighbour in $\Ev\setminus\{r\}$ needs an argument (an odd vertex with two neighbours
$a,a'\ne r$ is handled by Theorem~\ref{thm:pair} with any $u\in\Ev\setminus\{r,a,a'\}$). Let $D\ge2$
be the least distance between two distinct even vertices (finite, as $G$ is connected), attained by
a pair $a,t$ with $a\ne r$, and let $a=p_0,p_1,\dots,p_D=t$ be a shortest path. Its interior vertices
are odd, since an even one would be closer to $a$. Choose
$u\in\Ev\setminus\{r,a,t\}$, which exists since $e\ge4$. No $p_i$ ($0<i<D$) is adjacent to $u$:
otherwise $d(a,u)\le i+1$ and $d(u,t)\le D-i+1$, so minimality forces $D=2$, and then $p_1$ would
have the two neighbours $a,u$ in $\Ev\setminus\{r\}$. Now the sequence of moves
$(p_0,p_1),\dots,(p_{D-2},p_{D-1})$ (odd transfers; each $p_i$ is odd of degree at least $3$ and
not adjacent to $u$) followed by $(p_{D-1},t)$ (an even step; $t\notin\{p_{D-1},u\}$ and $t\not\sim
u$) is a Fan certificate for $\{a,u\}$, so $\Ev\setminus\{a,u\}$ is realisable by
Theorem~\ref{thm:fancert}.
\end{proof}

For $e=4$ a set $\Ev\setminus\{u,v\}$ with $u,v\ne r$ is a realisable set of size $e/2$ containing $r$,
so Theorem~\ref{thm:b} proves the case $e=4$ of Proposition~\ref{prop:small}, for all connected
graphs and without Botler--Sambinelli, except when $G[\Ev]=K_3+r$; Corollary~\ref{cor:four}
already does so with no exception. (The proof given for
Proposition~\ref{prop:small} yields a partition into $n/2$ paths with an end at $r$, which is (R),
but it does not show that this partition has a simple profile; Corollary~\ref{cor:R} allows an odd
vertex with three ends.) The exceptional configuration is not reached by these moves. In
the computations of Section~\ref{app:comp} the second step holds in every residual instance of
Corollary~\ref{cor:fd2strong} and Theorem~\ref{thm:b} through order $8$.


\subsection{Further reductions}

The following moves remove a vertex $x\in X$ by Lov\'asz's theorem directly. Each is stated for a
general $X\subseteq\Ev$; for $X=\{u,v\}$ it is a sufficient condition for the graph
$G^{\Ev\setminus\{u,v\}}$, or any graph with exactly two even vertices, to be tight. In each case
the smaller graph may have odd order, which is harmless by Remark~\ref{rem:anyorder}.

\begin{proposition}\label{prop:moves}
Let $X\subseteq\Ev$ and $x\in X$.
\begin{itemize}
\item[(L)] If $x$ has a neighbour $\ell$ of degree~$1$, and $\Ev\setminus X$ is realisable in
$G-\ell$, then it is realisable in $G$.
\item[(S)] If $d(x)=2$ and the neighbours $y,z$ of $x$ are not adjacent, and $\Ev\setminus X$ is
realisable in $G'=G-x+yz$, then it is realisable in $G$.
\item[(C)] If every neighbour of $x$ lies in $\Ev\setminus X$, and $\Ev\setminus X\setminus N(x)$ is
realisable in $G-x$, then $\Ev\setminus X$ is realisable in $G$.
\end{itemize}
In particular, if $X=\{x,v\}$ and one of the three local conditions holds at $x$, then
$\Ev\setminus X$ is realisable (in (C) the smaller set is $\Ev(G-x)\setminus\{v\}$, realisable by
Lemma~\ref{lem:minus}; in (L) and (S) the smaller graph has $\Ev\setminus\{x\}$ as its even set).
\end{proposition}

\begin{proof}
(L) In $G-\ell$ the vertex $x$ is odd and $\Ev(G-\ell)=\Ev\setminus\{x\}\supseteq\Ev\setminus X$; a
realisation $D'$ has one end at $x$; appending $x\ell$ to that path gives $x$ no end and $\ell$ one,
with $|D'|$ unchanged and $|\Od(G)|=|\Od(G-\ell)|$. (S) $G'$ is simple, $y$ and $z$ keep their
degrees, and $\Ev(G')=\Ev\setminus\{x\}$. In a realisation of $\Ev\setminus X$ in $G'$, replace the
edge $yz$ of its path by $y\,x\,z$; this is a path since $x\notin V(G')$, and $x$ is interior. (C) In
$G-x$ the neighbours of $x$ become odd, so
$\Ev(G-x)=\Ev\setminus(\{x\}\cup N(x))$; a realisation $D'$ of $\Ev\setminus X\setminus N(x)$ has one
end at each $y\in N(x)$. Pair $N(x)$ arbitrarily and add the $d(x)/2$ paths $y\,x\,y'$: each
$y\in N(x)$ now has two ends, $x$ has none, and
$|D'|+d(x)/2=(o+d(x))/2+|\Ev\setminus X|-d(x)+d(x)/2=o/2+|\Ev\setminus X|$.
\end{proof}

For a graph $H$ with exactly two even vertices $u,v$ (taking $G=H$, $X=\{u,v\}$) the moves give the
following sufficient conditions for tightness, which we state for later reference:
\begin{itemize}
\item $u$ has a pendant neighbour (L); or $d(u)=2$ with non-adjacent neighbours (S); or every
neighbour of $u$ is an odd vertex adjacent to a pendant vertex (delete one such pendant vertex at
each neighbour of $u$ to get $G_0$; then $H=G_0^{N(u)}$, and this is (C) for $G_0$ and $X=\{u,v\}$
through Lemma~\ref{lem:dict});
\item (\emph{transfer}) $u$ has an odd neighbour $w\not\sim v$ such that $H-uw$, whose even vertices
are $w$ and $v$, is tight (Lemma~\ref{lem:fantransfer} with $S=\emptyset$). Iterating: if a path
$u=x_0,x_1,\dots,x_k$ has $x_1,\dots,x_k$ odd and outside $N[v]$, and $x_k$ is adjacent to a
pendant vertex, then $H$ is tight; if $d(u)=2$, $N(u)=\{x,y\}$ and $x,y\notin N[v]$, then $H$ is tight
even when $xy\in E(H)$ (transfer along $u,x,y$, after which $u$ is a pendant vertex at $y$).
\end{itemize}

\subsection{Obstructions}\label{sec:obstructions}

Let $H$ have $n$ vertices, exactly two of them even, $u$ and $v$, so that the target is
$p=n/2-1$ paths.

\begin{proposition}[Obstructions]\label{prop:obstructions}
$H$ is not tight in each of the following cases.
\begin{itemize}
\item[(O1)] Some odd vertex has degree $n-1$.
\item[(O2)] $uv\in E(H)$, $d(u)=d(v)=2$, and $u,v$ have a common neighbour.
\item[(O3)] $uv\notin E(H)$, $N(u)=N(v)=\{x,y\}$, $xy\in E(H)$ and $d(y)=3$.
\end{itemize}
\end{proposition}

\begin{proof}
(O1) In any partition $D$, $d(w)+b_D(w)\le2|D|$, since a path has at most two edges at $w$ and one
fewer for each of its ends there. In a tight partition $|D|=|\Od|/2=n/2-1$ and $b_D(w)=1$ for odd
$w$, so $d(w)\le n-3$. (Formalized: \texttt{degree\_le\_of\_tight}.) (O2) With $u$ and $v$ interior,
the transitions at $u$ and at $v$ are forced and the triangle $u\,v\,x$ would lie on one path. (O3)
With $u,v$ interior the segments $x\,u\,y$ and $x\,v\,y$ lie on two paths, which are distinct (their
union is a cycle), and the edge $xy$ lies on a third path (a path containing $x\,u\,y$ cannot also
contain $xy$). All three paths contain $y$, which has only three edges, so $y$ is an end of all
three, whereas it must be an end of exactly one.
\end{proof}

(O1) contains the edge bound: if no odd vertex has degree $n-1$, then every odd vertex has at
least two non-neighbours and $u,v$ at least one each, so at least $n-1$ edges are missing and
$|E(H)|\le\binom n2-(n-1)=p(n-1)$, which is what $p$ paths of at most $n-1$ edges allow. The
graph $K_4-e$ of the introduction is (O1) (and (O3)).

\emph{Cut vertices.} Let $c$ be a cut vertex, $H=H_1\cup H_2$ with $V(H_1)\cap V(H_2)=\{c\}$ and
both sides non-trivial. Call a partition $D_i$ of $H_i$ \emph{admissible} if $b_{D_i}(x)$ equals
the required value in $H$ (one for odd $x$, zero for $x\in\{u,v\}$) at every $x\ne c$. Then $H$ is
tight if and only if there are admissible $D_1,D_2$ with
$|b_{D_1}(c)-b_{D_2}(c)|=\beta(c)$, where $\beta(c)\in\{0,1\}$ is the required value at $c$.
Indeed a path of $H$ through $c$ visits $c$ once, so cutting at $c$ the paths that use edges of
both sides gives admissible $D_1,D_2$ in which $k$ pairs of ends at $c$ come from the cut paths and
$b_{D_1}(c)+b_{D_2}(c)-2k=\beta(c)$; conversely two pieces ending at $c$ from different sides meet
only in $c$, so any $\min\{b_{D_1}(c),b_{D_2}(c)\}$ pairs of them can be glued. In particular $H$ is
not tight as soon as one side has no admissible partition at all; (O3) is the case
$H_1=H[\{x,y,u,v\}]\cong K_4-e$, $c=x$.

\emph{Census.} Among the connected graphs with exactly two even vertices there are $3$ of order
$4$ ($2$ not tight, both (O1)), none of order~$5$, $47$ of order $6$ ($17$ not tight: $16$ are (O1) and one
is (O2)), and $2668$ of order $8$, of which $372$ are not tight: $344$ are (O1), $21$ are (O2), $4$
are (O3), and the remaining $3$ have a cut vertex at which one side has no admissible partition
(a dense block next to a vertex with two pendant neighbours). In particular every non-tight
$2$-connected graph of order at most $8$ with exactly two even vertices is (O1). We ask whether
this holds in general:

\begin{conjecture}\label{conj:twoeven}
A $2$-connected graph with exactly two even vertices is tight unless some odd vertex is adjacent
to all other vertices.
\end{conjecture}

For the second descent step only graphs $G^{\Ev\setminus\{u,v\}}$ matter, in which every vertex other
than $u,v$ that is even in $G$ carries a pendant vertex; (O1) is then impossible once $e\ge4$, since
the order is $n+e-2$ and the degrees are at most $n$.

\subsection{What is left of the second step}

Theorems~\ref{thm:second}, \ref{thm:pair} and~\ref{thm:b} leave two configurations.
\begin{itemize}
\item \emph{With the first removed vertex $u$ prescribed} (Corollary~\ref{cor:fd2strong}(i)): the
even vertices are independent and the neighbourhoods $N(a)$, $a\in\Ev\setminus\{u,r\}$, are pairwise
disjoint sets of odd vertices. Fan certificates still apply when some $a\in\Ev\setminus\{u,r\}$ is
joined to an even vertex $t\ne u$ by a path whose interior vertices are odd and not adjacent to
$u$ (the proof of Theorem~\ref{thm:b}), and the moves (L), (S) apply at vertices of degree two or
with a pendant neighbour. What resists is a vertex $u$ whose neighbourhood meets every such path.
\item \emph{Four even vertices} (Theorem~\ref{thm:b}): $G[\Ev]=K_3+r$. The triangle vertices are
pairwise twins and $r$ sees none of them. For two triangle vertices $a,a'$ and the third one $u$,
an even common neighbour of $a,a'$ other than $u$ does not exist, and for an odd one the second Fan
step fails, since $aa'\in E$ and $u$, $t$ are both zero-end neighbours of $t'$. Certificates with
odd transfers still give $\Ev\setminus\{a,u\}$ for triangle vertices $a\ne u$ when $a$ has a
pendant neighbour (move (L)), or an odd neighbour $w$ of degree at least $3$ adjacent to neither
other triangle vertex (the moves $(a,w)$, then the even step $(u,b)$ to the third triangle vertex
$b$), or when $a$ is joined to $r$ by a path whose interior vertices are odd and not adjacent to
$u$. So the configuration is residual only if every odd neighbour of a triangle vertex sees at
least two triangle vertices and every path from a triangle vertex $a$ to $r$ meets the neighbourhoods of
both other triangle vertices in its interior (as in $K_6$ minus the three edges at $r$).
\end{itemize}
In both, the obstacle is the same: Fan's lemmas need fewer zero-end neighbours than available ends,
and the configuration places a removed vertex next to every candidate.

\begin{remark}[Both configurations are closed]\label{rem:twostep}
Theorem~\ref{thm:two} of Section~\ref{sec:four} proves the second step without exceptions:
\emph{if $G$ is connected of even order with $e\ge4$ and $y\ne w$ are even, then some even
$v\notin\{y,w\}$ has $\Ev\setminus\{w,v\}$ realisable.} It is an induction along a path from $y$ to
another even vertex, using the transfer across an even edge (Lemma~\ref{lem:edge}, which is the case
$S=\Ev\setminus\{v\}$ of Lemma~\ref{lem:fandescent}), the star lemma (Lemma~\ref{lem:star}) and the
lift along an odd neighbour (Lemma~\ref{lem:lift}). Taking $y=r$ and $w=u$ settles
Theorem~\ref{thm:b}, and Corollary~\ref{cor:fd2strong} for connected $G$, in full, including
$G[\Ev]=K_3+r$. Connectivity is needed: in $P_3$ together with a disjoint $K_5-e$, the root $r$ the
middle vertex of $P_3$, exception (ii) of Corollary~\ref{cor:fd2strong} occurs and no $v$ works,
since the $K_5-e$ component cannot be covered by two paths. The Lean
proof is \texttt{two\_step\_descent} in \texttt{lean/GallaiLean/EMCFour.lean} on the project's base
branch; we rebuilt it and checked that it uses only the standard axioms, with Lov\'asz as the hypothesis
\texttt{LovaszStatement}. The results of this section remain useful for their explicit witnesses:
the twin criterion of Corollary~\ref{cor:twins} says which pairs $\{u,v\}$ can fail, and
Theorem~\ref{thm:pair} names the pair among which $v$ can be chosen.
\end{remark}

\section{Stronger induction hypotheses, and a barrier for Fan moves}\label{sec:H}

Theorem~\ref{thm:branching} reduces Gallai's conjecture to the branching case $\mathrm{B}^0_k$, where
the even vertex $y$ has at least two even neighbours. There, certificates made of Lov\'asz's theorem
and Fan's transfers can stall even when (R) holds. Even an induction hypothesis of type~(R), with one
prescribed vertex, does not help: the realisation it supplies may have zero-end vertices next to every
re-inserted edge, and Lemma~\ref{lem:fan} then cannot fire. This section asks whether a stronger
hypothesis, one that controls where the seed has its zeros, repairs the transfer step. Throughout, $G$
has even order $n$ and no isolated vertex, $e=|\Ev|$, and $N_E(y)=N(y)\cap\Ev$.

The answer has three parts.
\begin{itemize}
\item One prescribed vertex suffices when $y$ has exactly two even neighbours. So the open case
becomes ``at least three even neighbours'' (Corollary~\ref{cor:twonbrs}).
\item Hypotheses with more prescribed vertices do not propagate. Computationally, the strongest
possible hypothesis on end profiles of $G-e$, together with Fan's moves on the single edge $e$,
certifies every even root of every connected graph of order at most $8$ except $K_8$ minus a perfect
matching (Remark~\ref{rem:oracle}).
\item That graph, and every $K_{2m}$ minus a perfect matching with $m\ge4$, cannot be finished by any
Fan move, whatever the seed (Theorem~\ref{thm:barrier}, Corollary~\ref{cor:cp}).
\end{itemize}

\begin{definition}[Prescribed descent]\label{def:prescribed}
For $P\subseteq\Ev$ and $k\ge0$, $G$ has \emph{$k$-step descent keeping $P$} if some $k$-set
$R\subseteq\Ev\setminus P$ has $\Ev\setminus R\in\cF(G)$. Since $\cF(G)$ is an up-set, this holds
exactly when some $S\in\cF(G)$ with $P\subseteq S$ has $|S|\le e-k$.
\end{definition}

For $P=\{y\}$ this is the notion in $\mathrm{D}^0_k$ and $\mathrm{B}^0_k$. At the critical level
$e=2k$, descent keeping $P$ says that $P$ lies in a basis of the endpoint matroid.

\begin{lemma}[Prescribed transfer]\label{lem:ptransfer}
Let $y,a$ be adjacent even vertices, $K=G-ya$, and $P\subseteq\Ev\setminus\{a\}$. If $K$ has $k$-step
descent keeping $(P\setminus\{y\})\cup(N_E(y)\setminus\{a\})$, then $G$ has $(k+1)$-step descent
keeping $P\cup\{y\}$. Conversely, if a single transfer of $ya$ towards $y$ turns a partition of $K$
that realises some $S$ into one that realises $S'$, then $\{y\}\cup N_E(y)\setminus\{a\}\subseteq S'$
and $a\notin S'$.
\end{lemma}

\begin{proof}
In $K$ the vertices $y,a$ are odd and $\Ev(K)=\Ev\setminus\{y,a\}$. The even neighbours of $y$ in $K$
are $N_E(y)\setminus\{a\}$; its other neighbours are odd. Let $R$ be the $k$-set of the hypothesis and
$S=\Ev(K)\setminus R$. Then $y$ has no neighbour in $\Ev(K)\setminus S$, so
Lemma~\ref{lem:transfer} gives $S\cup\{y\}=\Ev\setminus(R\cup\{a\})\in\cF(G)$. The set $R\cup\{a\}$ has
$k+1$ elements and avoids $P\cup\{y\}$. For the converse, Lemma~\ref{lem:fan} with $T=\{a\}$ needs
$|Z_D(y)|<b_D(a)=1$. By~\eqref{eq:passing} this means $N_E(y)\setminus\{a\}\subseteq S$. After the
transfer $y$ has two ends and $a$ none.
\end{proof}

\begin{corollary}[Two even neighbours reduce]\label{cor:twonbrs}
Let $k\ge1$ and assume $\mathrm{D}^0_k$. Let $G$ be connected with $e\ge 2k+2$. Let the even vertex
$y$ have exactly two even neighbours $a\ne b$, and let every bridge at $y$ be pendant. Then $G$ has
$(k+1)$-step descent keeping $y$. Hence $\mathrm{B}^0_k$ may be weakened to $\mathrm{B}^3_k$, where $y$
has at least three even neighbours. If $\mathrm{B}^3_k$ holds for every $k\ge3$, then Gallai's
conjecture holds.
\end{corollary}

\begin{proof}
The edge $ya$ is not a bridge, since $d(a)\ge2$. So $K=G-ya$ is connected of even order,
$\Ev(K)=\Ev\setminus\{y,a\}$ has at least $2k$ elements, and $b\in\Ev(K)$. By $\mathrm{D}^0_k$, $K$
has $k$-step descent keeping $b$. Since $N_E(y)\setminus\{a\}=\{b\}$, Lemma~\ref{lem:ptransfer} with
$P=\{y\}$ applies. In the induction of Theorem~\ref{thm:branching}, $\mathrm{D}^0_k$ is available at
level $k+1$, so $\mathrm{B}^3_{k+1}$ suffices there.
\end{proof}

{\sloppy In Lean (\texttt{lean/GallaiLean/FanBarrier.lean}) these results are \texttt{desc\_succ\_of\_transfer\_two},
\texttt{branchBase\_of\_three} and \texttt{gallai\_of\_branchBase3}: Gallai's conjecture follows from
the single named hypothesis \texttt{BranchBase3} for all $k\ge3$. Lemma~\ref{lem:ptransfer} is
formalised only in the case used by Corollary~\ref{cor:twonbrs} ($P=\{y\}$, two even neighbours);
its converse is not formalised.\par}

\begin{example}\label{ex:gqz}\sloppy
Let $G_q=K_8-\{03,12,23,34,56,67\}$ (graph6 code Gqz\textasciitilde vw). Then $\Ev=\{0,1,3,4,5,7\}$, $N(3)=\{1,5,6,7\}$ and
$N_E(3)=\{1,5,7\}$. The graph is critical, and all $20$ triples of even vertices are realisable, so its
endpoint matroid is uniform. Lov\'asz-seeded Fan certificates do not reach the root $3$ with up to
eight held edges. In the induction with (R) as hypothesis, the adversary places a zero at an even
neighbour of $3$.

Two prescribed vertices suffice. The graph $K=G_q-31$ has the four even vertices $0,4,5,7$, and all
six pairs are realisable in it. In particular $\{5,7\}=N_E(3)\setminus\{1\}$ is realisable, and
Lemma~\ref{lem:ptransfer} gives $\{3,5,7\}\in\cF(G_q)$. The edges $35$ and $37$ work in the same way.
So a hypothesis on $G_q-e$, for an edge $e$ at $3$, does yield a certificate at the root $3$. The
hypothesis needed is that a particular pair lies in the endpoint matroid of a four-even-vertex graph.
Theorem~\ref{thm:emc4} does not supply this: it allows $\{5,7\}$ to be one of the two disjoint
non-realisable pairs.
\end{example}

\begin{remark}[What prescribed descent propagates through]\label{rem:propagate}
Let $\mathrm{D}^{(j)}_k$ say that every connected graph of even order with $e\ge2k$ has $k$-step
descent keeping $P$, for every $P\subseteq\Ev$ with $|P|\le j$. Run the proof of
Theorem~\ref{thm:branching} at a vertex $y\in P$.
\begin{itemize}
\item \emph{Bridge $yt$.} The pieces need descent keeping $P\cap A$ and $P\cap B$, or
$(P\cap B)\cup\{t\}$ when $t$ is odd and $|A|$ is even. These sets have at most $j$ elements.
The exception is $t\in P$ with $t$ even and $|A|$ even: gluing~(2) of Lemma~\ref{lem:bridge} removes $t$.
\item \emph{Unique even neighbour $t$.} The step needs $\mathrm{D}^{(j)}_k$ in $G-yt$, and $t\notin P$.
\item \emph{Lift along an odd neighbour $t$.} The step needs $\mathrm{D}^{(j)}_{k+1}$ in $G-yt$
keeping $(P\setminus\{y\})\cup\{t\}$.
\item \emph{Branching with $m$ even neighbours.} Lemma~\ref{lem:ptransfer} needs
$\mathrm{D}^{(j+m-2)}_k$ in $G-ya$, with $a\notin P$.
\end{itemize}
So no bounded $j$ is closed under the reduction. For $j\ge2$ the statement is false at the critical
level: a critical graph with four even vertices and a non-realisable pair violates
$\mathrm{D}^{(2)}_2$, and such graphs exist at order $8$ (Section~\ref{app:comp}).

The matroid version asks that $P$ be independent in the endpoint matroid. Its branching step needs
some $a\in N_E(y)$ for which $N_E(y)\setminus\{a\}$ is independent in the endpoint matroid of
$G-ya$. Conjecture~\ref{conj:emc} for $G-ya$ says nothing about which sets are independent, and the
requirement is false in two ways.
\begin{itemize}
\item It fails whenever $m>e/2$. The set produced contains $\{y\}\cup N_E(y)\setminus\{a\}$.
\item It fails in small cases. Among the $275$ pairs $(G,y)$ with $G$ critical of order~$8$,
$e=6$ and $m=3$, there are $22$ in which no choice of $a$ works.
\end{itemize}
An example of the second kind is GsaEDo, with $y=7$ joined to a universal vertex by three ears
$7\,i\,0$, $i=2,3,4$. In these cases the pair $N_E(y)\setminus\{a\}$ is non-realisable in $G-ya$
for every choice of $a$.
\end{remark}

To measure how far any hypothesis on $G-F$ can go, we give the induction an oracle.

\begin{definition}[Oracle-seeded Fan game]\label{def:oracle}
Hold a nonempty set $F\subseteq E(G)$. As seed, take any realisable end profile of $K=G-F$, given by
a partition of $K$. In the \emph{simple} game the profile must realise a set $S\in\cF(K)$; in the
\emph{general} game it may be any profile. Then re-insert the edges of $F$ by the following moves:
\begin{itemize}
\item single transfers;
\item grouped transfers, where an adversary chooses the donor;
\item new one-edge paths.
\end{itemize}
The play wins for the root $r$ when $F$ is used up with $\sum b\le n$ and $b(r)\ge1$.
\end{definition}

A winning play gives a partition of $G$ into at most $n/2$ paths with an end at $r$, so it proves (R)
at $r$. The general game uses everything that any hypothesis about end profiles of $G-F$ could
provide.

\begin{remark}[The oracle game at order at most eight]\label{rem:oracle}
We solved the game exactly for every even root of every connected graph of order at most $8$
(programs \texttt{tools/realis.c} and \texttt{tools/fan\_oracle.py}). At order $8$ there are $11{,}117$
connected graphs and $44{,}008$ even roots. For comparison, Lov\'asz seeds with $|F|\le4$ lose on $50$
graphs of order $8$.

\begin{center}\small
\begin{tabular}{lp{0.5\textwidth}p{0.3\textwidth}}
order & simple game & general game \\ \hline
$4$ & $|F|=1$ wins everywhere & -- \\
$6$ & $|F|=1$: only $K_{2,2,2}$ loses (all roots); $|F|\le2$ wins everywhere & $|F|=1$ wins everywhere\\
$8$ & $|F|=1$: $6$ graphs lose; $|F|\le2$ and $|F|\le3$: $3$ graphs lose & $|F|=1$: only $M_8$ loses
\end{tabular}
\end{center}

Here $M_8=K_8$ minus a perfect matching.
\begin{itemize}
\item The three simple-game losers at $|F|\le3$ are $M_8$ (every root), Gqzn\textasciicircum[ (root $1$)
and Gut\textasciitilde\textasciitilde w (roots $1,3,4$). The last two win the general game with $|F|\le2$.
\item $M_8$ also loses the simple game with $|F|\le4$ and the general game with $|F|\le3$. By
Corollary~\ref{cor:cp} it loses for every $F$.
\end{itemize}
So at order at most $8$, knowing the realisable profiles of $G-e$ for one well-chosen edge $e$, and
playing Fan's moves on $e$, certifies (R) everywhere except on $M_8$. The endpoint matroid of $M_8$ is
uniform: all $70$ half-sets are realisable, so $M_8$ satisfies Conjecture~\ref{conj:emc} in its
strongest form.
\end{remark}

\begin{theorem}[The last-edge barrier]\label{thm:barrier}
Let $D$ be a partition of a graph $H$, and let $u\ne t$ be vertices with $ut\notin E(H)$. Write
$N^{\mathrm{ev}}_H(u)$ for the set of neighbours of $u$ of even degree in $H$.
\begin{enumerate}
\item If $b_D(t)>|Z_D(u)|$ (Fan's condition for a single transfer of $ut$ towards $u$), then
\[
b_D(u)+2\,|N^{\mathrm{ev}}_H(u)|+2\le 3|D|.
\]
\item If $H=G-F$ for a $d$-regular graph $G$ of order $n$, then $nd\le 2\bigl(|D|(n-1)+|F|\bigr)$.
\end{enumerate}
\end{theorem}

\begin{proof}
(1) Let $N=N^{\mathrm{ev}}_H(u)$ and $Z=Z_D(u)$. Every $w\in N\setminus Z$ has $b_D(w)\ge2$, because
$b_D(w)$ is positive and even by~\eqref{eq:parity}. Neither $u$ nor $t$ lies in $N$. By~\eqref{eq:parity},
\[
2|D|=\sum_v b_D(v)\ \ge\ b_D(u)+b_D(t)+2|N|-2|Z|.
\]
Now use $|Z|\le b_D(t)-1$, and $b_D(t)\le|D|$, which holds because a path has at most one end at $t$.
This gives $2|D|\ge b_D(u)+2|N|+2-b_D(t)\ge b_D(u)+2|N|+2-|D|$.
(2) Every path has at most $n-1$ edges, so $|E(G)|-|F|\le|E(H)|\le|D|(n-1)$. Also $2|E(G)|=nd$.
\end{proof}

\begin{corollary}[No Fan move finishes $K_{2m}$ minus a perfect matching]\label{cor:cp}
Let $m\ge4$, and let $\mathrm{CP}(m)$ be $K_{2m}$ minus a perfect matching.
\begin{enumerate}
\item For every edge $ut$ and every partition $D$ of $\mathrm{CP}(m)-ut$ that satisfies Fan's condition
for a transfer of $ut$ towards $u$, we have $|D|>m$.
\item If $D$ partitions $\mathrm{CP}(m)-F$ into at most $m-1$ paths, then $|F|\ge m-1$.
\end{enumerate}
Consider a play that ends with a partition of $\mathrm{CP}(m)$ into at most $m$ paths, built from a
partition of a proper spanning subgraph by Fan transfers and new paths made of held edges. Its last
move is a new path with at least $m-1$ edges, added to a partition of $\mathrm{CP}(m)-F$ into $m-1$
paths. In particular the last move is not a transfer, and it is not a new path of length one.
\end{corollary}

\begin{proof}
$\mathrm{CP}(m)$ is $(2m-2)$-regular. In $H=\mathrm{CP}(m)-ut$ the vertex $u$ is odd, so $b_D(u)\ge1$.
Its $2m-3$ neighbours have degree $2m-2$ in $H$. Theorem~\ref{thm:barrier}(1) gives
$3|D|\ge 1+2(2m-3)+2=4m-3>3m$. Part~(2) of the theorem with $|D|\le m-1$ gives
\[
2m(2m-2)\le 2\bigl((m-1)(2m-1)+|F|\bigr),\quad\text{that is,}\quad |F|\ge m-1.
\]
A grouped transfer re-inserts one of its held edges and leaves the others held. So the last transfer
of a play is a single transfer into $H=\mathrm{CP}(m)-ut$, and a transfer preserves $|D|$. By~(1)
the last move is therefore a new path. It consists of all the held edges $F$ that remain, and it is
added to a partition $D$ of $\mathrm{CP}(m)-F$ with $|D|\le m-1$, so $|F|\ge m-1$ by~(2). In fact
$|D|=m-1$: with $|D|\le m-2$, part~(2) of the theorem would give
$|F|\ge 2m(m-1)-(m-2)(2m-1)=3m-2$, more than the $2m-1$ edges of a path in $\mathrm{CP}(m)$.
\end{proof}

{\sloppy Corollary~\ref{cor:cp} is formalised as \texttt{cocktail\_fan\_barrier} and
\texttt{cocktail\_insertion\_barrier}, and Theorem~\ref{thm:barrier} as \texttt{fan\_transfer\_barrier} and
\texttt{insertion\_barrier\_regular}. All are in \texttt{lean/GallaiLean/FanBarrier.lean}, use no
hypotheses, and depend only on the standard axioms. Part~(1) of the theorem is general: for any
even-regular graph with $d>(3n-2)/4$, no Fan transfer can finish at the budget $n/2$.\par}

\emph{Credit.} For $K_{2m}$ minus a perfect matching, the failure of Fan's numerical test is due
to~\cite[Theorem~10.21]{SP}. It is proved there in the rooted form $(F-e)^t$, for both orientations
of every edge and every partition within budget. Corollary~\ref{cor:cp}(1) re-derives it in the
leafless form. Theorem~\ref{thm:barrier}(1) is the general counting behind it.

\begin{remark}[Consequence for stronger hypotheses]\label{rem:Hverdict}
An induction hypothesis about realisable sets or end profiles of proper spanning subgraphs enters a
Lov\'asz--Fan reduction only through the seed. In the reduction of Theorem~\ref{thm:branching}, every
step other than bridge gluing is a Fan transfer: the transfer of Lemma~\ref{lem:ptransfer}, and the lift
of Theorem~\ref{thm:reduction}. Bridge gluing does not apply to the $(2m-2)$-connected graphs
$\mathrm{CP}(m)$. By Corollary~\ref{cor:cp}, no such hypothesis, however strong, makes the reduction
reach $\mathrm{CP}(m)$ for $m\ge4$. Yet $\mathrm{CP}(4)=M_8$ is critical and its endpoint matroid is
uniform. The obstruction lies in the moves, not in the hypothesis.

A proof must spend a spare path: a new path of at least $m-1$ edges, added to a partition of
$\mathrm{CP}(m)-F$ into $m-1$ paths. For $m=4$ and a new path of exactly $m-1=3$ edges, the remaining $21$
edges form three Hamiltonian paths. This is exactly the flower configuration by which the companion manuscript~\cite{MM}
(Proposition~5.5) proves (R) on $M_8$. For every $m$ it is the explicit construction
of~\cite[Theorem~10.21]{SP}. Its last move adds the transversal path $L$ of three edges to
the three paths $C_i-e_i$, so the bound $|F|\ge m-1$ is attained. The two results together locate the
missing ingredient. It is not a stronger hypothesis; it is a move that spends a spare path at
Fan's equality.
\end{remark}

\subsection*{The flower move}

Corollary~\ref{cor:cp} says how a certificate for $\mathrm{CP}(m)$ must end: with a new path of at
least $m-1$ edges, added to a partition of the rest into $m-1$ paths. For $\mathrm{CP}(m)$ this ending
is known to exist~\cite[Theorem~10.21]{SP}. Our contributions here are the following:
\begin{itemize}
\item the flower lemma in the generality and formalisation of Lemma~\ref{lem:flower} (the lemma itself
appears in~\cite{MM} and in the proof of~\cite[Theorem~38]{MR});
\item the fact that Fan's equality forces a flower (Proposition~\ref{prop:equality});
\item the receipt-or-flower dichotomy (Theorem~\ref{thm:receiptflower});
\item its application to $B'$ (Proposition~\ref{prop:Bprime}).
\end{itemize}
We now show that the ending arises from Fan's receipts. It comes from the state in which Fan's receipts at a hub stop at equality. Throughout, a
path $P$ from $a$ to $b$ with $ab\notin E(P)$ is \emph{closed} into the cycle $C(P)=P+ab$.

\begin{lemma}[Flowers; {\cite[Lemma~5.4]{MM}}, {\cite[proof of Theorem~38]{MR}}]\label{lem:flower}
Let $P_1,\dots,P_k$ be paths, and let $L$ be a path whose edges are $e_1,\dots,e_k$ with $e_i\in
E(C(P_i))$. Then the paths $C(P_i)-e_i$ ($i\le k$) and $L$ have, together, exactly the edges of the
cycles $C(P_i)$, there are $k+1$ of them, and they have two ends at every vertex of $L$ and none
elsewhere.
\end{lemma}

\begin{proof}
$C(P_i)-e_i$ is a path whose ends are the two ends of $e_i$. So a vertex $v$ is an end once for each
$e_i$ containing $v$, and once for each end of $L$ at $v$. This gives two at an interior vertex of
$L$ (two edges of $L$), two at an end of $L$ (one edge and one end), and none elsewhere.
\end{proof}

\begin{proposition}[Equality forces a flower]\label{prop:equality}
Let $D$ be a partition of $H$, let $s$ be a vertex and $T$ a set of vertices with $s\notin T$, and
suppose that $b_D(s)=|T|$, $b_D=1$ on $T$ and $b_D=0$ elsewhere. Then:
\begin{itemize}
\item $|D|=|T|$;
\item every path of $D$ joins $s$ to a vertex of $T$;
\item every $t\in T$ is the end of exactly one path $P_t$.
\end{itemize}
If $st\notin E(H)$ for $t\in T$, the cycles $C(P_t)$ are edge-disjoint, pass through $s$, and cover
$E(H)\cup\{st:t\in T\}$.
\end{proposition}

\begin{proof}
By~\eqref{eq:parity}, $2|D|=\sum b_D=2|T|$. A path has at most one end at $s$, and $b_D(s)=|D|$, so
every path has exactly one end at $s$. All other ends lie in $T$, one at each vertex.
\end{proof}

This is the state in which the receipts at a hub end.
\begin{itemize}
\item Let $K$ be $G$ with every edge at $s$ removed, and let $D_0$ be a partition of $K$ with one
end at every neighbour of $s$ and none elsewhere. Then $|D_0|=d(s)/2$. Suppose $d(s)=n-2$, every
neighbour of $s$ is even, and the one non-neighbour $s^*\ne s$ is even. Then $s^*$ is the only even
vertex of $K$, and Lov\'asz's theorem gives such a $D_0$.
\item Re-insert the edges $sy$ by Fan's receipts. After $j$ receipts, $s$ has $j$ ends and $j$ blank
neighbours, the donors. The $d(s)-j$ unreceived neighbours have one end each.
\item The next grouped receipt has supply $d(s)-j$ against $j$ blanks. So it succeeds as long as
$j<d(s)/2$.
\item At $j=d(s)/2$ the supply equals the blank count, Fan's lemma stops, and
Proposition~\ref{prop:equality} applies with $T$ the set of unreceived neighbours.
\end{itemize}
Either a further receipt succeeds, or the paths close into cycles through $s$.

\begin{theorem}[The flower move]\label{thm:flowermove}
In the situation of Proposition~\ref{prop:equality}, suppose that some path $L$ through $s$ has
exactly one edge from each cycle $C(P_t)$, a \emph{transversal}. Then $E(H)\cup\{st:t\in T\}$ has a
partition into $|D|+1$ paths with two ends at every vertex of $L$ and none elsewhere. In
particular $b(s)=2$.
\end{theorem}

\begin{proof}
Lemma~\ref{lem:flower}, with Proposition~\ref{prop:equality}.
\end{proof}

The move spends one spare path, so it stays within the budget $n/2$ exactly when $|D|\le n/2-1$.
For a Lov\'asz seed of $G$ minus the star at $s$, this holds whenever $d(s)\le n-2$. No edge count
enters the move itself. The only condition is that a transversal exists, and $L$ has $|T|$ edges.
For $\mathrm{CP}(m)$ this is $m-1$, the least length that Corollary~\ref{cor:cp} allows.

\begin{theorem}[Cocktail party graphs]\label{thm:cpm}
Let $m\ge4$. Suppose that every system of $m-1$ edge-disjoint Hamiltonian cycles through a vertex
$s$, on $2m$ vertices, has a transversal through $s$. Then $\mathrm{CP}(m)$ satisfies the rooted
test at every vertex.
\end{theorem}

\begin{proof}
The proof has five steps.
\begin{itemize}
\item \emph{Seed.} In $K=\mathrm{CP}(m)$ minus the star at $s$, the $2m-2$ neighbours of $s$ have
degree $2m-3$, and the antipode $s^*$ is the only even vertex. Lov\'asz gives $m-1$ paths with one
end at each neighbour of $s$.
\item \emph{Receipts.} After $j$ receipts the next one has supply $2m-2-j$ against $j$ blank
neighbours. This exceeds $j$ for $j\le m-2$, so $m-1$ receipts succeed. The $m$-th meets equality.
\item \emph{Cycles.} By Proposition~\ref{prop:equality}, the $m-1$ paths close into edge-disjoint
cycles through $s$ that cover $E(\mathrm{CP}(m))$.
\item \emph{Hamiltonicity.} The cycles have $2m(m-1)=(m-1)\cdot 2m$ edges in total, and each has at
most $2m$. So all of them are Hamiltonian.
\item \emph{Flower.} A transversal through $s$ and Theorem~\ref{thm:flowermove} give $m$ paths with
two ends at $s$.
\end{itemize}
\end{proof}

\begin{proposition}[Transversals for at most four cycles]\label{prop:rt4}
Let $C_1,\dots,C_k$ be edge-disjoint Hamiltonian cycles on a common vertex set, with $k\le4$. Then
every vertex $s$ lies on a path with exactly one edge from each $C_i$.
\end{proposition}

\begin{proof}
An edge lies on only one of the cycles, and every vertex has exactly two neighbours on each.

\emph{The greedy path.} Choose a $C_1$-neighbour $a_1$ of $s$ and a $C_2$-neighbour $a_2$ of $a_1$.
Then $a_2\ne s$, since $sa_1\in C_1$. The two $C_3$-neighbours of $a_2$ do not include $a_1$, since
$a_1a_2\in C_2$, and at most one of them is $s$. So some $C_3$-neighbour $a_3$ of $a_2$ lies outside
$\{s,a_1\}$. The path $s\,a_1\,a_2\,a_3$ settles $k\le3$.

\emph{The rotation, for $k=4$.} If a $C_4$-neighbour $y$ of $a_3$ lies outside $\{s,a_1\}$, take
$s\,a_1\,a_2\,a_3\,y$; note $y\ne a_2$, since $a_2a_3\in C_3$. Otherwise the two $C_4$-neighbours of
$a_3$ are $s$ and $a_1$. Let $x\ne a_2$ be the other $C_2$-neighbour of $a_1$. Then $x\ne s$, since
$sa_1\in C_1$, and $x\ne a_3$, since $a_1a_3\in C_4$. The path $a_2\,a_3\,s\,a_1\,x$ uses the edges
$a_2a_3\in C_3$, $a_3s\in C_4$, $sa_1\in C_1$ and $a_1x\in C_2$.
\end{proof}

\begin{corollary}\label{cor:cp45}
$\mathrm{CP}(4)=M_8$ and $\mathrm{CP}(5)$ satisfy the rooted test at every vertex.
\end{corollary}

\emph{Credit and status.} The rooted test for $\mathrm{CP}(m)$ at every vertex and every $m\ge4$
is Theorem~10.21 of~\cite{SP}. Its proof is a flower with an explicit transversal:
\begin{itemize}
\item take Walecki's Hamilton cycles on $\{t\}\cup\mathbb Z_{2h}$;
\item detour each antipodal missing edge through the mate of $t$, which gives $h=m-1$ Hamilton cycles;
\item open all of them along the path $Q=(t,1,\dots,h)$, which meets them in turn, and adjoin $Q$.
\end{itemize}
So Theorem~\ref{thm:cpm} and Corollary~\ref{cor:cp45} are not new results about $\mathrm{CP}(m)$.
They show only that the same ending is reached from an \emph{arbitrary} Lov\'asz seed through Fan's
receipts. That route needs a transversal of whatever Hamilton cycles the receipts produce.
\begin{itemize}
\item Proposition~\ref{prop:rt4} provides one for at most four cycles.
\item Marked-reduction Theorem~38~\cite{MR} gives more without any transversal: $j\le7$
edge-disjoint cycles through a vertex $x$, Hamiltonian or not, partition into at most $j+1$ paths
with at least two ends at $x$. The path count is Chu, Fan and Zhou's Lemma~2.5. So the receipt route
already closes $\mathrm{CP}(m)$ for $m\le8$.
\item Theorem~13.18 of~\cite{SP} gives neighbour-exchange conditions for opening any number of
cycles through a vertex.
\end{itemize}
For $\mathrm{CP}(4)$ the receipt route is Proposition~5.5 of~\cite{MM}.

\begin{question}[Rainbow transversals]\label{qu:rt}
Let $C_1,\dots,C_k$ be edge-disjoint Hamiltonian cycles on $2k+2$ vertices. Does every vertex $s$
lie on (or even start) a path with exactly one edge from each $C_i$?
\end{question}

For $\mathrm{CP}(m)$ the question is moot: \cite[Theorem~10.21]{SP} settles the graphs directly. We
keep it only as an independent combinatorial question about arbitrary Hamilton decompositions, and
about the receipt route of Theorem~\ref{thm:cpm}.

\emph{What is known.} Question~\ref{qu:rt} asks for a rainbow path of a given length in an
edge-colouring. Here every vertex meets each of the $k$ colours exactly twice, the path must have
length $k$ (so it uses every colour), and it must pass through a prescribed vertex.
\begin{itemize}
\item \emph{Arbitrary colour degree.} Let every vertex meet at least $k$ colours. Chen and
Li~\cite{CL05,CL06} prove a rainbow path of length at least $k-1$ for $3\le k\le7$, and at least
$\lceil 2k/3\rceil+1$ for $k\ge8$, which confirms a conjecture of Saito. They conjecture $k-1$ in
general. Length $k$ is not available at this generality.
\item \emph{Proper colourings of $K_n$.} Andersen~\cite{And89} conjectured a rainbow path of length
$n-2$. Asymptotic versions are known~\cite{APS17}. Gyárfás and Mhalla~\cite{GM10} prove length at
least $(2n+1)/3$ for $1$-factorisations. Maamoun and Meyniel~\cite{MaMe84} give $1$-factorisations of
$K_{2^t}$ without a rainbow Hamilton path.
\item \emph{Graph systems.} Transversal Hamilton paths and cycles in a system of graphs are known
under Dirac-type degree conditions (Joos and Kim~\cite{JK20}, and the survey literature since).
Our graphs have degree $2$, so these results do not apply.
\end{itemize}
Our case lies outside all of these. The length equals the colour degree, which exceeds the Chen--Li
bound, and it is attainable only because each colour class is $2$-regular on at least $2k+1$
vertices. We found no result that covers it.

\emph{Evidence and partial results.}
\begin{itemize}
\item Proposition~\ref{prop:rt4} settles $k\le4$. The proof uses only that every vertex has
exactly two neighbours of each colour, not that the classes are Hamiltonian. The missing-move
manuscript~\cite[Lemma~5.6]{MM} gives another proof for $k\le4$.
\item The stronger form (a transversal \emph{starting} at $s$) held at every vertex in every test
for $m\le10$ (\texttt{tools/rainbow\_transversal.py}). The tests used $60$ random Hamiltonian
decompositions of $\mathrm{CP}(m)$ and a Walecki-type one per $m$.
\item The same holds for random $2$-factorisations of $\mathrm{CP}(m)$, $m\le8$, in which the colour
classes have short cycles (\texttt{tools/rainbow\_variants.py}). So Hamiltonicity does not seem to
be the point, and a local argument may exist.
\item Neither random greedy nor a Warnsdorff-type greedy succeeds without backtracking. The random
greedy fails $5$--$15\%$ of the time. The Warnsdorff rule fails in about $1\%$ of cases from
$m=8$. A proof therefore needs rotations, as in Proposition~\ref{prop:rt4}.
\item A relaxation is far easier in practice. Call a path a \emph{segment transversal} if it is a
concatenation of one contiguous segment of each cycle, each with at least one edge. Lemma~\ref{lem:flower} extends to segments verbatim:
$C_i$ minus a segment is a path, and $L$ has two ends at its segment junctions and at its ends. A
no-backtracking rule found a segment transversal starting at $s$ in every test ($m\le11$).
\end{itemize}

\subsection*{Receipt or flower: the move in general}

The hub flower applies beyond $\mathrm{CP}(m)$. The rigid alternative of~\cite[Lemma~5.1]{MM} is
exactly its hypothesis.

Marked-reduction Corollary~39~\cite{MR} has the same structure in a different setting. A new
vertex is joined to ends of old paths, and the doubly assigned paths close through it into cycles,
which Theorem~38 of~\cite{MR} then opens. The next theorem is the version at Fan's equality in the
receipt process.

\begin{theorem}[Receipt or flower]\label{thm:receiptflower}
Let $D$ be a partition of $H$, and let $s$ be a vertex with $b_D(s)=d_H(s)$. Let $T$ be a set of
non-neighbours of $s$, with $s\notin T$ and $b_D=1$ on $T$, and put $G=H+\{st:t\in T\}$. Then one of the following
holds.
\begin{enumerate}
\item Some $t\in T$ ends a path of $D$ that avoids $s$. Appending $ts$ to that path gives a partition
of $H+st$ into $|D|$ paths, with one more end at $s$ and none at $t$.
\item Every path ending in $T$ joins $T$ to $s$. These $|T|$ paths close into edge-disjoint cycles
$C_t$ through $s$. For every transversal $L$ of the cycles $C_t$, the graph $G$ has a partition into
$|D|+1$ paths, namely the cut cycles, $L$ and the other paths of $D$, with
\[
b'(v)=b_D(v)-|T|\,[v=s]-[v\in T]+2\,[v\in V(L)].
\]
\end{enumerate}
\end{theorem}

\begin{proof}
Suppose (1) fails. Then every path ending at some $t\in T$ contains $s$. Since $b_D(s)=d_H(s)$, no
path passes through $s$, so every path containing $s$ ends there. Thus each path ending in $T$
joins $T$ to $s$. These paths are distinct, because each has only one end in $T$. Lemma~\ref{lem:flower}
applied to them, with the remaining paths left unchanged, gives~(2).
\end{proof}

{\sloppy In Lean the recombination step of~(2) is \texttt{flower\_decomp\_rest}. Case~(2) is useful
when an edge count makes the paths $C_t-st$ Hamiltonian on a common vertex set, as in
Theorem~\ref{thm:cpm}. Then Proposition~\ref{prop:rt4} supplies a transversal when $|T|\le4$. With a
Lov\'asz seed of $G$ minus the star at $s$, and receipts from all of $N(s)$, case~(2) leaves no
other paths exactly when the seed has ends only on $N(s)$. With at most one even vertex in the seed
graph, this means: \emph{$G$ is Eulerian, of even order $n$, and $d(s)=n-2$.} The move then spends
the spare path and meets the budget $n/2$ exactly.\par}

\begin{proposition}[The graph $B'$]\label{prop:Bprime}
Let $B'$ have a clique $U$ on six vertices, a vertex $x_1$ adjacent to $U$ and to a leaf $\ell$, and
two adjacent vertices $x_3,x_4$ adjacent to all of $U$. This is the residual case Q2 of the $d=6$
frontier~\cite[\S5]{MM}, with $n=10$. Then (R) holds at every vertex of $B'$, by a Lov\'asz seed
of $B'-x_1$, three receipts at $x_1$, and Theorem~\ref{thm:receiptflower}.
\end{proposition}

\begin{proof}
In $K=B'$ with the star at $x_1$ removed, the vertices $x_1$ and $\ell$ are isolated. All other
vertices have degree $7$. By Lov\'asz, $K$ has a partition into $4$ paths with one end at each
vertex of $U\cup\{x_3,x_4\}$. Three receipts at $x_1$ from $U$ have supplies $1>0$, $5>1$, $4>2$.
They give a partition $D$ of $H$, the graph $K$ plus three spokes, with $b(x_1)=3=d_H(x_1)$, $b=1$
on the set $T$ of the three unreceived core vertices and on $x_3,x_4$, and $|E(H)|=31$. Apply
Theorem~\ref{thm:receiptflower}.

\emph{Case (1).} Append $tx_1$. Then append the leaf edge $x_1\ell$ to a path ending at $x_1$, and
add the new path $t'x_1t''$ on the other two spokes. The result has $5$ paths, and $b=2$ at the
core vertices $t',t''$.

\emph{Case (2).} The three paths through $x_1$ end in $T$. The fourth path $R$ has the remaining
ends $x_3,x_4$. All four avoid $\ell$, and $R$ also avoids $x_1$, since $b(x_1)=d_H(x_1)$. So the
paths have at most $8,8,8$ and $7$ edges. As $31=3\cdot8+7$, all bounds are attained, and the
cycles $C_t$ are Hamiltonian on $V\setminus\{\ell\}$. Proposition~\ref{prop:rt4} gives a transversal
$L=x_1a_1a_2a_3$ with $a_1\in U$. Theorem~\ref{thm:receiptflower} gives $4+1=5$ paths, and
appending $x_1\ell$ at the end $x_1$ of $L$ keeps $5$ paths, with $b(a_1)=2$.

In both cases some core vertex is an end of a partition into $n/2$ paths. Every permutation of $U$
is an automorphism of $B'$, so this holds at every core vertex. At the odd vertices it holds by
parity.
\end{proof}

The same reading applies to $A_{6,2}$ (\cite[Proposition~5.2]{MM}). There the rigid case is
excluded by the count $30>3\cdot8$, so case (1) occurs. The resulting partition into $5$ paths has
$b=2$ at the two core vertices $t',t''$ of the new path $t'st''$. By the symmetry of the core, this
gives (R) at every core vertex of $A_{6,2}$ as well. These are the core roots that an earlier draft
of~\cite{MM} listed as open; its Proposition~5.2 now includes them.

{\sloppy In Lean (\texttt{lean/GallaiLean/Flower.lean}), Lemma~\ref{lem:flower} is
\texttt{flower\_decomp}, built from \texttt{exists\_close\_cut}, and Proposition~\ref{prop:equality}
is \texttt{hub\_cycles}. Theorem~\ref{thm:cpm} is \texttt{cocktail\_passes\_of\_rt} and
\texttt{cocktail\_rooted\_of\_rt}, with the hypothesis \texttt{RainbowTransversal}.
Proposition~\ref{prop:rt4} is \texttt{rainbowTransversal\_three} and
\texttt{rainbowTransversal\_four} in the cases needed there, $k=3$ on $8$ vertices and $k=4$ on
$10$ vertices, and Corollary~\ref{cor:cp45} is \texttt{cocktail\_passes\_four} and
\texttt{cocktail\_passes\_five}; the general recombination of Theorem~\ref{thm:receiptflower}(2) is
\texttt{flower\_decomp\_rest}. There is no \texttt{sorry}, Lov\'asz's theorem is the library's
proof, and the axioms are the standard ones.\par}

\section{Computational evidence}\label{app:comp}

The accompanying programs (\texttt{tools/} in the project repository) compute exact partitions by
enumerating simple paths and solving an exact-cover integer program (OR-Tools CP-SAT). Every
partition returned was independently checked. On all connected graphs of order at most six,
Conjectures~\ref{conj:emc} and~\ref{conj:fd} hold; the realisable family of every critical graph of
order four or six is the spanning family of a loopless matroid, uniform whenever the graph is
$2$-connected. At order eight the same holds for all $1{,}082$ connected critical graphs: the
minimal realisable sets are the bases of a loopless matroid of rank $e/2$. The matroid is uniform for
$1{,}030$ of them, including all $476$ three-connected ones. The $52$ non-uniform ones have a cut
vertex ($46$) or connectivity two ($6$, among them GsaFC\{ of Proposition~\ref{prop:uhc}). This run
uses the exact depth-first oracle \texttt{tools/realis.c}, with its own pruning. The oracle was checked
against the integer program on a sample of graphs of orders six and eight (\texttt{tools/emc\_sweep.py}). For four even vertices, Corollary~\ref{cor:four} was checked
on all connected graphs of order at most six (critical or not) and on the first 500 of the $5{,}523$
connected graphs of order eight with four even vertices, with no exception (\texttt{tools/emc\_four.py});
the run was stopped once the corollary was proved. In every critical graph found with a non-realisable
pair there are exactly two such pairs, and they are disjoint. The examples of Propositions~\ref{prop:cp}--\ref{prop:uhc}
were found by the same exhaustive searches.

For Section~\ref{sec:second} (program \texttt{tools/two\_even.py}): the census of graphs with
exactly two even vertices quoted in Section~\ref{sec:obstructions} is exhaustive through order~$8$.
Through order $6$, every non-realisable pair $\{u,v\}$ is an E-twin pair ($384$ pairs), the pair
step has no counterexample ($565$ instances) and every set certified by a Fan schedule is
realisable ($657$ sets); these are checks of the proofs, which are also formalized. Through
order~$6$ the second step of faithful descent succeeds in every residual instance of
Corollary~\ref{cor:fd2strong} ($9$ instances) and Theorem~\ref{thm:b} ($3$ instances); at order
$8$ there are $1357$ and $399$ such instances, and the step succeeds in all of them.
Conjecture~\ref{conj:twoeven} also holds for $85$ random $2$-connected graphs of orders $10$ and
$12$ with exactly two even vertices and no universal odd vertex (program
\texttt{tools/two\_even\_random.py}).

\begin{thebibliography}{9}
\bibitem{BS} F.~Botler and M.~Sambinelli, Towards Gallai's path decomposition conjecture, J.~Graph Theory 97 (2021) 161--184.
\bibitem{Fan} G.~Fan, Path decompositions and Gallai's conjecture, J.~Combin.\ Theory Ser.~B 93 (2005)
117--125.
\bibitem{Cai} M.~Cai, Faithful path partitions (as cited in the short-faithful manuscript).
\bibitem{L68} L.~Lov\'asz, On covering of graphs, in: Theory of Graphs (Proc.\ Colloq.\ Tihany, 1966),
Academic Press, 1968, 231--236.
\bibitem{MR} J.~Williams, Cut vertices and cycle intersections in Gallai's path decomposition problem,
manuscript, November 2026.
\bibitem{REC} J.~Williams, Root endpoints and cut inequalities in path decompositions, manuscript,
November 2026.
\bibitem{And89} L.~D.~Andersen, Hamilton circuits with many colours in properly edge-coloured
complete graphs, Math.\ Scand.\ 64 (1989) 5--14.
\bibitem{APS17} N.~Alon, A.~Pokrovskiy and B.~Sudakov, Random subgraphs of properly edge-coloured
complete graphs and long rainbow cycles, Israel J.\ Math.\ 222 (2017) 317--331 (arXiv:1608.07028).
\bibitem{CL05} H.~Chen and X.~Li, Long heterochromatic paths in edge-colored graphs, Electron.\ J.\
Combin.\ 12 (2005) R33.
\bibitem{CL06} H.~Chen and X.~Li, Color degree and color neighborhood union conditions for long
heterochromatic paths in edge-colored graphs, arXiv:math/0512144.
\bibitem{GM10} A.~Gyárfás and M.~Mhalla, Rainbow and orthogonal paths in factorizations of $K_n$,
J.\ Combin.\ Des.\ 18 (2010) 167--176.
\bibitem{JK20} F.~Joos and J.~Kim, On a rainbow version of Dirac's theorem, Bull.\ Lond.\ Math.\ Soc.\
52 (2020) 498--504 (arXiv:1910.01281).
\bibitem{MaMe84} M.~Maamoun and H.~Meyniel, On a problem of G.~Hahn about coloured Hamiltonian paths
in $K_{2^t}$, Discrete Math.\ 51 (1984) 213--214.
\bibitem{MM} J.~Williams, Lov\'asz's spare path: the missing move in the Fan game, manuscript,
\texttt{manuscripts/missing-move} of the project repository.
\bibitem{SP} J.~Williams, Separators and endpoint exchange in the rooted Gallai problem, manuscript,
November 2026.
\end{thebibliography}
\end{document}
